% ══════════════════════════════════════════════════════════════════
% THESIS MAIN — From Reactive Detection to Proactive Consensus
% A Byzantine Fault Tolerant Blockchain Security Evaluation Framework
% for Advanced Metering Infrastructure
%
% VJTI Mumbai — Department of Electrical Engineering
% Under the guidance of Prof. Uday Suryavanshi \& Dr. Sunny Kumar
% 16-Chapter Revised Structure
% ══════════════════════════════════════════════════════════════════

\documentclass[12pt, oneside, a4paper]{report}

% ─── Core Packages ───────────────────────────────────────────────────
\usepackage[utf8]{inputenc}
\usepackage[T1]{fontenc}
\usepackage{lmodern}
\usepackage{tcolorbox}
\tcbuselibrary{skins,breakable}
\usepackage[labelfont=bf, font=small, labelsep=period]{caption}

% ─── Mathematics ─────────────────────────────────────────────────────
\usepackage{amsmath, amssymb, amsthm, mathtools}
\usepackage{bm}
\usepackage{dsfont}

% ─── Page Layout ─────────────────────────────────────────────────────
\usepackage[
  a4paper,
  top=2.5cm,
  bottom=2.5cm,
  left=3.5cm,
  right=2.5cm
]{geometry}
\usepackage{setspace}
\onehalfspacing

% ─── Headers and Footers ─────────────────────────────────────────────
\usepackage{fancyhdr}
\pagestyle{fancy}
\fancyhf{}
\fancyhead[L]{\nouppercase{\leftmark}}
\fancyhead[R]{\thepage}
\renewcommand{\headrulewidth}{0.4pt}
\fancypagestyle{plain}{%
  \fancyhf{}
  \fancyhead[R]{\thepage}
  \renewcommand{\headrulewidth}{0pt}
}

% ─── Chapter and Section Formatting ──────────────────────────────────
\usepackage{titlesec}
\titleformat{\chapter}[display]
  {\normalfont\Huge\bfseries\sffamily\raggedleft}
  {\MakeUppercase{\chaptertitlename}%
    \rlap{ \resizebox{!}{1.5cm}{\thechapter} \rule{5cm}{1.5cm}}}
  {10pt}{\Huge}
\titlespacing*{\chapter}{0pt}{30pt}{40pt}

% ─── Tables ──────────────────────────────────────────────────────────
\usepackage{booktabs, multirow, array, longtable, tabularx, colortbl}
\usepackage{rotating}

% ─── Figures ─────────────────────────────────────────────────────────
\usepackage{graphicx, subcaption, float}
\graphicspath{{../figures/}{figures/}{../../figures/}}

% ─── Algorithms ──────────────────────────────────────────────────────
\usepackage{algorithm}
\usepackage{algorithmic}

% ─── Code Listings ───────────────────────────────────────────────────
\usepackage{listings}
\usepackage{xcolor}
\definecolor{codegreen}{rgb}{0,0.6,0}
\definecolor{codegray}{rgb}{0.5,0.5,0.5}
\definecolor{codeblue}{rgb}{0,0,0.8}
\definecolor{backcolour}{rgb}{0.97,0.97,0.97}
\lstdefinestyle{pythonstyle}{
  backgroundcolor=\color{backcolour},
  commentstyle=\color{codegreen},
  keywordstyle=\color{codeblue},
  numberstyle=\tiny\color{codegray},
  stringstyle=\color{red!70!black},
  basicstyle=\ttfamily\footnotesize,
  breakatwhitespace=false,
  breaklines=true,
  keepspaces=true,
  numbers=left,
  numbersep=5pt,
  showspaces=false,
  showstringspaces=false,
  showtabs=false,
  tabsize=2,
  language=Python
}

% ─── Colours ─────────────────────────────────────────────────────────
\definecolor{ieee}{RGB}{0,83,159}

% ─── Hyperlinks ──────────────────────────────────────────────────────
\usepackage[
  colorlinks=true,
  linkcolor=ieee,
  citecolor=ieee,
  urlcolor=ieee,
  bookmarks=true,
  pdfauthor={Atharv Manojkumar Shukla, VJTI Mumbai},
  pdftitle={From Reactive Detection to Proactive Consensus},
  pdfsubject={BFT Blockchain Security for Advanced Metering Infrastructure}
]{hyperref}

% ─── Cross-references ────────────────────────────────────────────────
\usepackage[noabbrev, capitalise]{cleveref}

% ─── Bibliography ────────────────────────────────────────────────────
\usepackage[square, numbers, sort&compress]{natbib}


% ─── Custom Environments ─────────────────────────────────────────────
\newtcolorbox{premiumbox}[1][]{
  enhanced,
  boxrule=0.5pt,
  colback=blue!5!white,
  colframe=ieee,
  arc=4pt,
  drop fuzzy shadow,
  #1
}

% ─── Utilities ───────────────────────────────────────────────────────
\usepackage{enumitem}
\usepackage{microtype}
\usepackage{tocloft}
\usepackage{emptypage}

% ─── Theorem Environments ────────────────────────────────────────────
\newtheorem{theorem}{Theorem}[chapter]
\newtheorem{lemma}[theorem]{Lemma}
\newtheorem{proposition}[theorem]{Proposition}
\newtheorem{corollary}[theorem]{Corollary}
\theoremstyle{definition}
\newtheorem{definition}[theorem]{Definition}
\newtheorem{assumption}[theorem]{Assumption}
\theoremstyle{remark}
\newtheorem{remark}[theorem]{Remark}
\newtheorem{example}[theorem]{Example}

% ─── Custom Commands ─────────────────────────────────────────────────
\newcommand{\analytical}{[\textit{Analytical}]}
\newcommand{\simulated}{[\textit{Simulated}]}
\newcommand{\litderived}{[\textit{Literature-Derived}]}
% Probability shorthand
\newcommand{\Pbyz}{P_{\mathrm{Byz}}}
\newcommand{\Pids}{P_{\mathrm{IDS}}}
\newcommand{\Pcomp}{P_{\mathrm{comp}}}
\newcommand{\Pvuln}{P_{\mathrm{vuln}}}
\newcommand{\Psa}{P_{\mathrm{SA}}}
\newcommand{\Pca}{P_{\mathrm{CA}}}
\newcommand{\Pscada}{P_{\mathrm{SCADA}}}
\renewcommand{\Pr}{P_{\mathrm{R}}}
\newcommand{\Wvuln}{W_{\mathrm{vuln}}}

% ─── List of Algorithms ──────────────────────────────────────────────
\renewcommand{\listalgorithmname}{List of Algorithms}

% ══════════════════════════════════════════════════════════════════════
\begin{document}

% ────────────────────────────────────────────────────────────────
% FRONT MATTER
% ────────────────────────────────────────────────────────────────
\pagenumbering{roman}


% --- BEGIN INPUT: frontmatter/cover ---
% ══════════════════════════════════════════════════════════════════
% COVER PAGE
% ══════════════════════════════════════════════════════════════════
\begin{titlepage}
\centering

\vspace*{1.5cm}

% Institution Logo Placeholder
% \includegraphics[width=3.5cm]{../figures/vjti_logo.png}\\[0.8cm]

{\large \textbf{Veermata Jijabai Technological Institute}}\\[0.2cm]
{\large Department of Electrical Engineering}\\[0.2cm]
{\large Mumbai --- 400 019, India}\\[2.5cm]

\rule{\linewidth}{0.5mm}\\[0.4cm]
{\LARGE \bfseries From Reactive Detection to Proactive Consensus:}\\[0.3cm]
{\Large \bfseries A Comparative Evaluation of Byzantine Fault Tolerant}\\[0.3cm]
{\Large \bfseries Blockchain Security for Advanced Metering Infrastructure}\\[0.4cm]
\rule{\linewidth}{0.5mm}\\[1.5cm]

{\large \emph{A Research Project Report submitted in partial fulfilment}}\\[0.1cm]
{\large \emph{of the requirements of the}}\\[0.3cm]
{\large \textbf{Bachelor of Technology in Electrical Engineering}}\\[2.0cm]

\begin{minipage}[t]{0.45\textwidth}
  \begin{flushleft}
    \textbf{Submitted by:}\\[0.2cm]
    Atharv Manojkumar Shukla\\
    (3rd year B.Tech Undergrad at IIIT Nagpur)\\[0.5cm]
    \textbf{Project Guides:}\\[0.2cm]
    Prof.\ Uday Suryavanshi \& Dr.\ Sunny Kumar\\
    E-MC\textsuperscript{2} Laboratory\\
    Veermata Jijabai Technological Institute, Mumbai - Electrical Engineering Department
  \end{flushleft}
\end{minipage}
\hfill
\begin{minipage}[t]{0.45\textwidth}
  \begin{flushright}
    \textbf{Academic Year:}\\[0.2cm]
    2025--2026\\[0.5cm]
    \textbf{Submitted on:}\\[0.2cm]
    July 2026\\[0.5cm]
    \textbf{Comparative Extension of:}\\[0.2cm]
    \emph{Faquir et al.\ (SEGAN, 2022)}\\
    \emph{Sheikh et al.\ (IEEE Access, 2020)}
  \end{flushright}
\end{minipage}

\vfill

{\small The work presented in this report has been carried out at the E-MC\textsuperscript{2} Laboratory, Electrical Engineering Department, Veermata Jijabai Technological Institute, Mumbai, under the supervision of Prof.\ Uday Suryavanshi \& Dr.\ Sunny Kumar.}

\end{titlepage}

% --- END INPUT: frontmatter/cover ---

\cleardoublepage


% --- BEGIN INPUT: frontmatter/declaration ---
% ══════════════════════════════════════════════════════════════════
% DECLARATION
% ══════════════════════════════════════════════════════════════════
\chapter*{Declaration}
\addcontentsline{toc}{chapter}{Declaration}

I hereby declare that the work presented in this thesis titled \emph{``From Reactive Detection to Proactive Consensus: A Comparative Evaluation of Byzantine Fault Tolerant Blockchain Security for Advanced Metering Infrastructure''} is a faithful record of original research carried out by me under the supervision of Prof.\ Uday Suryavanshi \& Dr.\ Sunny Kumar, E-MC\textsuperscript{2} Laboratory, Veermata Jijabai Technological Institute, Mumbai - Electrical Engineering Department.

\medskip

This work has not been submitted, either in whole or in part, for any other degree or examination at any other university or institution. All sources of information used in this thesis have been duly acknowledged in the references.

\medskip

I further declare that the probabilistic models, simulation frameworks, and experimental analyses presented herein are my original contributions unless explicitly stated otherwise through appropriate literature citations. The comparative analyses extend and build upon the published works of Sheikh~\emph{et al.}~\cite{sheikh2020} and Faquir~\emph{et al.}~\cite{ids_smart_meter}, as clearly indicated throughout.

\bigskip\bigskip

\begin{tabular}{ll}
  \textbf{Signature:} & \underline{\hspace{5cm}} \\[1.0cm]
  \textbf{Name:}      & Atharv Manojkumar Shukla \\[0.4cm]
  \textbf{Affiliation:}& 3rd year B.Tech Undergraduate at IIIT Nagpur \\[0.4cm]
  \textbf{Date:}      & \today \\[0.4cm]
  \textbf{Place:}     & Mumbai, India \\
\end{tabular}

\bigskip\bigskip

\noindent\textbf{Supervised by:}

\bigskip

\begin{tabular}{ll}
  \textbf{Signature:} & \underline{\hspace{5cm}} \\[1.0cm]
  \textbf{Names:}     & Prof.\ Uday Suryavanshi \& Dr.\ Sunny Kumar \\[0.4cm]
  \textbf{Laboratory:}& E-MC\textsuperscript{2} Laboratory \\[0.4cm]
  \textbf{Department:}& Veermata Jijabai Technological Institute, Mumbai - Electrical Engineering Department \\
\end{tabular}

% --- END INPUT: frontmatter/declaration ---

\cleardoublepage


% --- BEGIN INPUT: frontmatter/dedication ---
% ══════════════════════════════════════════════════════════════════
% DEDICATION
% ══════════════════════════════════════════════════════════════════
\thispagestyle{empty}
\vspace*{4cm}

\begin{center}
  \itshape
  \large

  \emph{To my parents, whose sacrifice made this possible.}

  \vspace{0.5cm}

  \emph{To Prof.\ Uday Suryavanshi and Dr.\ Sunny Kumar,}\\
  \emph{whose guidance transformed a project into a discipline.}

  \vspace{0.5cm}

  \emph{And to the engineers of VJTI who built before me,}\\
  \emph{so that I may build further.}
\end{center}

% --- END INPUT: frontmatter/dedication ---

\cleardoublepage


% --- BEGIN INPUT: frontmatter/acknowledgements ---
% ══════════════════════════════════════════════════════════════════
% ACKNOWLEDGEMENTS
% ══════════════════════════════════════════════════════════════════
\chapter*{Acknowledgements}
\addcontentsline{toc}{chapter}{Acknowledgements}

This thesis is the product of sustained guidance, encouragement, and intellectual challenge from a number of individuals to whom I owe a sincere debt of gratitude.

\medskip

First and foremost, I extend my deepest thanks to my project guide, \textbf{Prof. Uday Suryavanshi \& Dr. Sunny Kumar}, Department of Electrical Engineering, Veermata Jijabai Technological Institute, Mumbai. His exacting standards for analytical rigour, his insistence that every claim be traceable to a derivation, and his ability to distinguish between genuine contribution and academic fluff have fundamentally shaped both this thesis and my understanding of what research means. The framework in this work grew directly from hours of discussion, critique, and iterative revision under his direction.

\medskip

I am grateful to the \textbf{Department of Electrical Engineering, VJTI Mumbai}, for providing the academic environment and computational resources that made this research possible.

\medskip

I would like to acknowledge the authors of the two foundational papers that form the comparative basis of this thesis: \textbf{Faquir~\emph{et al.}}~\cite{ids_smart_meter}, whose two-stage intrusion detection system for smart meters represents the reactive defence paradigm evaluated herein; and \textbf{Sheikh~\emph{et al.}}~\cite{sheikh2020}, whose static probabilistic security model for Byzantine consensus on the IEEE 33-bus system forms the mathematical cornerstone upon which our 12-attack extension is built.

\medskip

I also acknowledge the broader community of researchers whose published works are cited extensively throughout this thesis. In particular, the foundational contributions of \textbf{Lamport, Shostak, and Pease}~\cite{lamport1982} on the Byzantine Generals Problem, and \textbf{Castro and Liskov}~\cite{castro1999} on Practical Byzantine Fault Tolerance, defined the theoretical landscape within which all the consensus protocols evaluated in this thesis operate.

\medskip

On a personal note, I thank my family for their unconditional support, and my colleagues at VJTI for countless discussions that sharpened my arguments and revealed their weaknesses.

\medskip

Finally, I acknowledge that all shortcomings in this work remain entirely my own responsibility.

\bigskip

\begin{flushright}
  \textit{Atharv Manojkumar Shukla}\\
  \textit{VJTI Mumbai, July 2026}
\end{flushright}

% --- END INPUT: frontmatter/acknowledgements ---

\cleardoublepage


% --- BEGIN INPUT: frontmatter/abstract ---
% ══════════════════════════════════════════════════════════════════
% ABSTRACT (~1200 words)
% ══════════════════════════════════════════════════════════════════
\chapter*{Abstract}
\addcontentsline{toc}{chapter}{Abstract}

\noindent\textbf{Background and Motivation.}
The proliferation of Advanced Metering Infrastructure (AMI) across modern smart grid deployments has introduced an unprecedented scale of networked cyber-physical devices into energy distribution systems. A contemporary AMI deployment may encompass tens of thousands of smart meters, each connected to head-end systems (HES), meter data management systems (MDMS), and SCADA control centres through heterogeneous communication networks spanning home-area networks (HAN), neighbourhood-area networks (NAN), and wide-area networks (WAN). While this connectivity enables real-time demand response, energy theft detection, and dynamic pricing, it simultaneously exposes the grid to a growing and sophisticated portfolio of cyber-physical attacks, including false data injection, denial of service, man-in-the-middle interception, Sybil node impersonation, and Byzantine validator compromise.

\medskip

\noindent\textbf{Research Gap.}
The prevailing security paradigm in AMI deployments is \emph{reactive}: anomalies are detected after they have manifested in measured data streams, and a response is then issued. Two-stage intrusion detection systems (IDS), such as that proposed by Faquir~\emph{et al.}~\cite{ids_smart_meter}, represent the state of the art in this paradigm, combining support-vector machine (SVM) anomaly classification with temporal fault propagation graph (TFPG) pattern matching to detect known attack signatures. Despite their detection accuracy, such systems share a fundamental architectural limitation: they assume that the underlying data transport layer is honest. A compromised validator, a replayed measurement, or a coordinated Sybil cluster can systematically undermine the IDS's evidence base before any detection algorithm is invoked. The IDS detects consequences; it cannot prevent the consensus layer from accepting fraudulent data in the first instance.

\medskip

\noindent\textbf{Proposed Framework.}
This thesis proposes a \emph{Static-Temporal Security Framework} (STSF) that provides a principled analytical basis for comparing the reactive IDS paradigm against a proactive Byzantine Fault Tolerant (BFT) blockchain consensus paradigm across twelve categorised cyber-physical attack vectors. The framework extends the static four-component probabilistic model of Sheikh~\emph{et al.}~\cite{sheikh2020} — originally calibrated to the IEEE 33-bus distribution system with 51 nodes and 3,854 sensors — to a twelve-attack parallel failure model in which the joint system compromise probability is computed under Binomial Byzantine threshold analysis and Shamir Secret Sharing assumptions. The STSF further introduces a temporal dimension via a Poisson arrival process model that quantifies the vulnerability window of each security paradigm as a function of the attack arrival rate $\lambda$ and the detection/consensus latency $\tau$.

\medskip

\noindent\textbf{Mathematical Formulation.}
Under the proposed probabilistic model, the static system compromise probability for a BFT consensus layer is expressed as the cumulative Binomial distribution $P_{\mathrm{Byz}}(f,n) = \sum_{k=f+1}^{n} \binom{n}{k} p_c^k (1-p_c)^{n-k}$, where $n$ is the validator set size, $f = \lfloor(n-1)/3\rfloor$ is the Byzantine fault tolerance threshold, and $p_c$ is the per-node compromise probability. Across the twelve attack categories, the joint IDS survival probability is modelled as $P_{\mathrm{IDS}} = \prod_{j=1}^{12}(1 - P_j \cdot \mathds{1}[\text{attack}_j \text{ undetected}])$, reflecting the parallel failure model in which the IDS's defence collapses if any one category evades detection. The temporal framework defines the paradigm-specific \emph{vulnerability window} $W(\tau) = 1 - e^{-\lambda \tau}$ as the probability that at least one attack arrives during the response latency interval.

\medskip

\noindent\textbf{Protocol Evaluation.}
Nine BFT consensus protocols are evaluated within the framework, spanning four evolutionary generations: \textit{G1 Classical BFT} (Oral Message OM(m), Classic PBFT); \textit{G2 Committee-Delegated BFT} (IBFT~2.0, QBFT); \textit{G3 Hierarchical PBFT} (CE-PBFT, G-PBFT, SV-PBFT); and \textit{G4 Sub-Second BFT} (Tower~BFT, RVR). Each protocol is analysed for message complexity, consensus latency, Byzantine fault tolerance bound, and security feature coverage (threshold signatures, verifiable random functions, validator rotation, credit scoring). The analysis is calibrated against published Hyperledger Fabric benchmarks and the IEEE 33-bus system parameters from Sheikh~\emph{et al.}

\medskip

\noindent\textbf{Implementation.}
The STSF is implemented as a modular Python analytical engine comprising three components: (i) a pure-analytical probabilistic model that reproduces the Sheikh~\emph{et al.}\ baseline exactly before extending it; (ii) a Monte Carlo validation engine using $10^5$ trials per configuration for statistical confirmation; and (iii) a figure-generation pipeline producing publication-ready plots for all experimental comparisons. All source code is documented in Appendix~B.

\medskip

\noindent\textbf{Principal Findings.}
Under the assumptions of the presented framework, the following principal findings are obtained. First, within the static model, the proactive BFT consensus layer reduces the joint system compromise probability by approximately 170 orders of magnitude relative to the reactive IDS baseline across the twelve attack vectors. Second, within the temporal model, the vulnerability window of reactive IDS systems is found to be between $3.7\times$ and $31\times$ wider than that of the fastest evaluated BFT protocols, depending on the detection latency assumed. Third, among the evaluated protocols, \emph{Tower~BFT} and \emph{RVR} achieve the strongest combined security profile, with consensus latencies of 242.9~ms and 200~ms respectively, compared to 7,650~ms for Classic~PBFT. Fourth, the \emph{single point of failure} (SPOF) risk inherent in centralised IDS architectures is formally quantified and shown to be absent in the $f$-fault-tolerant consensus design.

\medskip

\noindent\textbf{Contributions.}
The principal novel contributions of this thesis are: (1) the first systematic comparative framework for IDS and BFT blockchain security applied to AMI under twelve attack categories; (2) a Poisson-process temporal security model that extends Sheikh~\emph{et al.}'s static framework to time-varying threat environments; (3) a four-generation evolutionary taxonomy of BFT protocols with AMI-specific feature scoring; (4) formal derivations of Binomial Byzantine threshold, Shamir sharing, and parallel failure models in the AMI context; and (5) a fully reproducible Python analytical engine calibrated to the IEEE 33-bus system.

\medskip

\noindent\textbf{Limitations and Future Work.}
All quantitative findings are derived from the analytical and simulation framework and are calibrated to published parameters; they are not based on physical deployment measurements. Future work includes deployment on a physical Hyperledger Fabric test network, incorporation of federated learning within the IDS baseline for stronger comparison, and extension to 5G-connected smart meter topologies.

\bigskip

\noindent\textbf{Keywords:} Advanced Metering Infrastructure; Byzantine Fault Tolerance; Blockchain Consensus; Intrusion Detection Systems; Smart Grid Security; PBFT; Tower BFT; RVR; Static-Temporal Security Framework; IEEE 33-bus system.

% --- END INPUT: frontmatter/abstract ---

\cleardoublepage


% --- BEGIN INPUT: frontmatter/roadmap ---
% ══════════════════════════════════════════════════════════════════
% RESEARCH ROADMAP & CONTRIBUTION MAPPING
% ══════════════════════════════════════════════════════════════════
\chapter*{Research Roadmap}
\addcontentsline{toc}{chapter}{Research Roadmap}

\noindent
To orient the reader through this investigation, this page provides a visual flowchart of the research stages and a complete mapping of our core research contributions to their respective implementation chapters, empirical experiments, and generated figures and tables.

\vspace{1.5cm}

\subsection*{Research Flowchart}

\begin{center}
\setlength{\tabcolsep}{12pt}
\renewcommand{\arraystretch}{1.2}
\begin{tabular}{c}
  \framebox[14cm]{\textbf{Phase I: Problem Definition \& Literature Synthesis (Chapters 1--5)}} \\
  $\downarrow$ \\
  \framebox[14cm]{\textbf{Phase II: Threat \& Attack Surface Taxonomy (Chapter 6)}} \\
  $\downarrow$ \\
  \framebox[14cm]{\textbf{Phase III: Analytical Foundation (Chapters 7--8)}} \\
  \begin{tabular}{cc}
    \framebox[6.5cm]{Security Math (Chapter 7)} & \framebox[6.5cm]{Consensus Math (Chapter 8)} \\
  \end{tabular} \\
  $\downarrow$ \\
  \framebox[14cm]{\textbf{Phase IV: Simulation Framework \& Evaluation Platform (Chapter 9)}} \\
  $\downarrow$ \\
  \framebox[14cm]{\textbf{Phase V: Methodology, Verification \& Validation (Chapters 10--12)}} \\
  $\downarrow$ \\
  \framebox[14cm]{\textbf{Phase VI: Empirical Analysis, Discussion \& Impact (Chapters 13--16)}} \\
\end{tabular}
\end{center}

\vspace{1.5cm}

\subsection*{Contribution Mapping Matrix}

The matrix below traces the key contributions of this thesis through the research artifacts, confirming complete coverage and scientific traceability.

\vspace{0.5cm}

\begin{table}[H]
\centering
\small
\renewcommand{\arraystretch}{1.3}
\begin{tabularx}{\textwidth}{lcp{3.5cm}cc}
\toprule
\textbf{Contribution} & \textbf{Chapter} & \textbf{Core Artifact / Concept} & \textbf{Experiment} & \textbf{Key Fig. / Tab.} \\
\midrule
\textbf{C1: Paradigm Framework} & 2, 7 & Comparison of reactive IDS vs.\ proactive BFT & E1, E2 & Fig.~13.2, Tab.~7.3 \\
\textbf{C2: Threat Expansion}   & 6, 7 & 12-attack model & E2, E7 & Fig.~13.3, Tab.~6.1 \\
\textbf{C3: Temporal Extension} & 7, 8 & Poisson-process temporal framework & E4, E10 & Fig.~13.5, Tab.~7.4 \\
\textbf{C4: Protocol Taxonomy}  & 5, 8 & Four-generation evolutionary model & E5, E9 & Fig.~13.6, Tab.~8.1 \\
\textbf{C5: Open Science Code}  & 9, 12 & Python engine & E6 & Fig.~13.7, Listing~B.1 \\
\bottomrule
\end{tabularx}
\caption{Traceability mapping matrix linking core thesis contributions to chapters, experiments, and visual assets.}
\label{tab:contribution_mapping}
\end{table}

% --- END INPUT: frontmatter/roadmap ---

\cleardoublepage

\tableofcontents
\cleardoublepage

\listoffigures
\cleardoublepage

\listoftables
\cleardoublepage

\listofalgorithms
\addcontentsline{toc}{chapter}{\listalgorithmname}
\cleardoublepage


% --- BEGIN INPUT: frontmatter/symbols ---
% ══════════════════════════════════════════════════════════════════
% LIST OF SYMBOLS
% ══════════════════════════════════════════════════════════════════
\chapter*{List of Symbols}
\addcontentsline{toc}{chapter}{List of Symbols}

Below is a summary of the principal mathematical symbols and variables used in the mathematical security and consensus formulations throughout this thesis.

\vspace{1.0cm}

\begin{longtable}{lll}
\toprule
\textbf{Symbol} & \textbf{Description} & \textbf{Unit / Domain} \\
\midrule
$x$ & System-wide vulnerability parameter & $[0, 1]$ \\
$N_{\mathrm{sen}}$ & Total number of sensors in the AMI grid & Integer ($3{,}854$) \\
$K_1$ & Total number of sensor pairs & Integer ($7{,}426{,}681$) \\
$K_2$ & Total number of network segments & Integer ($1{,}271$) \\
$n$ & Total number of consensus validator nodes & Integer \\
$f$ & Byzantine fault tolerance threshold limit & Integer ($\lfloor (n-1)/3 \rfloor$) \\
$p_c$ & Per-node validator compromise probability & $[0, 1]$ \\
$\lambda$ & Cyber attack arrival rate parameter & Attacks / second \\
$\tau$ & Security system response or consensus latency & Seconds \\
$T$ & Mission operating time & Seconds \\
$\Psa$ & Pairwise sensor compromise probability & $[0, 1]$ \\
$\Pca$ & Communication channel compromise probability & $[0, 1]$ \\
$\Pscada$ & SCADA database compromise probability & $[0, 1]$ \\
$\Pr$ & Receiver terminal compromise probability & $[0, 1]$ \\
$\Pids$ & Joint static reactive IDS compromise probability & $[0, 1]$ \\
$\Pbyz$ & Byzantine consensus failure probability & $[0, 1]$ \\
$\Wvuln$ & Vulnerability exposure window probability & $[0, 1]$ \\
$\eta$ & Key compromise threshold mitigation factor & $[0, 1]$ \\
$\beta$ & Byzantine fault tolerance mitigation factor & $[0, 1]$ \\
$\omega_{\mathrm{credit}}$ & Credit-weighted Sybil mitigation factor & $[0, 1]$ \\
$\mathrm{STSI}$ & Static-Temporal Security Index value & $[0, 1]$ \\
\bottomrule
\end{longtable}

% --- END INPUT: frontmatter/symbols ---

\cleardoublepage


% --- BEGIN INPUT: frontmatter/abbreviations ---
% ══════════════════════════════════════════════════════════════════
% LIST OF ABBREVIATIONS
% ══════════════════════════════════════════════════════════════════
\chapter*{List of Abbreviations}
\addcontentsline{toc}{chapter}{List of Abbreviations}

Below is a summary of the principal acronyms and abbreviations used throughout this thesis.

\vspace{1.0cm}

\begin{longtable}{ll}
\toprule
\textbf{Abbreviation} & \textbf{Full Description} \\
\midrule
ACL & Access Control List \\
AMI & Advanced Metering Infrastructure \\
APT & Advanced Persistent Threat \\
BDD & Bad Data Detection \\
BFT & Byzantine Fault Tolerance \\
BLS & Boneh-Lynn-Shacham (Signatures) \\
CAP & Consistency, Availability, Partition Tolerance \\
CTMC & Continuous-Time Markov Chain \\
DDoS & Distributed Denial of Service \\
DN & Distribution Network \\
DNS & Domain Name System \\
DoS & Denial of Service \\
ECDSA & Elliptic Curve Digital Signature Algorithm \\
EV & Electric Vehicle \\
FDI & False Data Injection \\
G1--G4 & Generation 1 to Generation 4 (Consensus) \\
GST & Global Stabilization Time \\
HAN & Home Area Network \\
HES & Head-End System \\
HIL & Hardware-in-the-Loop \\
HSM & Hardware Security Module \\
IBFT & Istanbul Byzantine Fault Tolerance \\
IDS & Intrusion Detection System \\
IoT & Internet of Things \\
JTAG & Joint Test Action Group \\
MDMS & Meter Data Management System \\
MMPP & Markov-Modulated Poisson Process \\
NAN & Neighborhood Area Network \\
NS-3 & Network Simulator 3 \\
PBFT & Practical Byzantine Fault Tolerance \\
PKI & Public Key Infrastructure \\
PLC & Power Line Communication \\
PoH & Proof-of-History \\
PoS & Proof-of-Stake \\
PoW & Proof-of-Work \\
QBFT & Quorum Byzantine Fault Tolerance \\
RBAC & Role-Based Access Control \\
RMSE & Root Mean Square Error \\
RTU & Remote Terminal Unit \\
RVR & Random Verifiable Rotation \\
SCADA & Supervisory Control and Data Acquisition \\
SPOF & Single Point of Failure \\
SSS & Shamir Secret Sharing \\
STSF & Static-Temporal Security Framework \\
STSI & Static-Temporal Security Index \\
TFPG & Temporal Fault Propagation Graph \\
TLS & Transport Layer Security \\
TPS & Transactions Per Second \\
V\&V & Verification and Validation \\
V2G & Vehicle-to-Grid \\
VRF & Verifiable Random Function \\
WAN & Wide Area Network \\
\bottomrule
\end{longtable}

% --- END INPUT: frontmatter/abbreviations ---

\cleardoublepage

% ────────────────────────────────────────────────────────────────
% MAIN MATTER
% ────────────────────────────────────────────────────────────────
\pagenumbering{arabic}


% --- BEGIN INPUT: chapters/ch01_introduction ---
% Chapter 1: Introduction


% --- END INPUT: chapters/ch01_introduction ---

\cleardoublepage


% --- BEGIN INPUT: chapters/ch02_background ---
% Chapter 2: Background


% --- END INPUT: chapters/ch02_background ---

\cleardoublepage


% --- BEGIN INPUT: chapters/ch03_threat_model ---
% ══════════════════════════════════════════════════════════════════
% CHAPTER 6: THREAT MODEL
% ══════════════════════════════════════════════════════════════════
\chapter{Introduction}
\label{ch:introduction}
% TODO: Introduction content

\chapter{Research Framework}
\label{ch:research_framework}
% TODO: Research Framework content

\chapter{AMI Background}
\label{ch:ami_background}
% TODO: AMI Background content

\chapter{IDS Background}
\label{ch:ids_background}
% TODO: IDS Background content

\chapter{Blockchain Background}
\label{ch:blockchain_background}
% TODO: Blockchain Background content

\chapter{Threat Model}
\label{ch:threat_model}

This chapter details the cyber-physical threat landscape of Advanced Metering Infrastructure (AMI). We systematically evaluate twelve distinct cyber-physical attack vectors. Every attack vector is presented under a uniform architectural template:
\begin{enumerate}[label=(\roman*)]
  \item \textbf{Architecture:} Physical and network configuration of the attack target.
  \item \textbf{Attack Steps:} Procedural sequence of the intrusion.
  \item \textbf{Impact:} Resulting damage to grid operations or market settlement.
  \item \textbf{Existing Mitigation:} How reactive IDS systems detect the intrusion.
  \item \textbf{IDS Failure Mode:} Why reactive IDS fails to provide complete prevention.
  \item \textbf{Blockchain Mitigation:} How distributed proactive BFT consensus blocks the attack.
  \item \textbf{Mathematical Probability Model:} Corresponding variables and probability terms.
  \item \textbf{Simulation Parameters:} Baseline parameters mapped in our evaluation model.
\end{enumerate}

All details in this chapter represent analytical threat analysis and literature-derived security mappings \analytical.

\bigskip

% ─────────────────────────────────────────────────────────────────
\section{Sensor Compromise}
\label{sec:tm:sensor_compromise}
% ─────────────────────────────────────────────────────────────────

\subsection{Architecture}
The sensor compromise attack targets the physical sensor endpoints located on household meters, charging stations, and grid buses (transformers).

\subsection{Attack Steps}
\begin{enumerate}
  \item The adversary acquires physical access to the device or exploits firmware vulnerabilities via regional network sweeps.
  \item The firmware is replaced with compromised binaries, or debug interfaces (JTAG) are accessed.
  \item The attacker modifies measurement data before it undergoes local encryption or signing.
\end{enumerate}

\subsection{Impact}
Compromised sensors send inaccurate voltage, current, and energy consumption values, leading to incorrect state estimation at the control center.

\subsection{Existing Mitigation}
Reactive IDS systems analyze historical load shapes and log-likelihood ratios to identify statistical anomalies in reading patterns.

\subsection{IDS Failure Mode}
If the attacker introduces a slow drift in readings rather than abrupt spikes, the anomaly remains within the standard variance limits, evading detection.

\subsection{Blockchain Mitigation}
By storing sensor measurements on a distributed ledger with signature verification, the data is secured against modification once written. Coordinated validator checks ensure that forged data is rejected during consensus before it can be appended to the state.

\subsection{Mathematical Probability Model}
The pairwise sensor compromise probability $\Psa(x)$ is:
\begin{equation}
\Psa(x) = x^2 (1 - x)^{N_{\mathrm{sen}} - 2}
\end{equation}

\subsection{Simulation Parameters}
The baseline probability is $P_1 = \Psa(0.5) \approx 1.6 \times 10^{-3}$, with sensor count $N_{\mathrm{sen}} = 3{,}854$.

\bigskip

% ─────────────────────────────────────────────────────────────────
\section{False Data Injection (FDI)}
\label{sec:tm:fdi}
% ─────────────────────────────────────────────────────────────────

\subsection{Architecture}
FDI targets the state estimation algorithms at the control center, manipulating telemetry measurements transmitted from Remote Terminal Units (RTUs) and smart meters.

\subsection{Attack Steps}
\begin{enumerate}
  \item The adversary analyzes the system topology matrix $H$ through leakage or public documents.
  \item An attack vector $a$ is constructed such that $a = H c$, where $c$ represents the target injection state.
  \item Telemetry vector $z$ is modified to $z' = z + a$ during communication transit.
\end{enumerate}

\subsection{Impact}
State estimators calculate clean measurement residuals:
\begin{equation}
r = \| z' - H \hat{x}' \| = \| z + H c - H(\hat{x} + c) \| = \| z - H \hat{x} \|
\end{equation}
The residual checks pass, allowing bad data to alter control center actions.

\subsection{Existing Mitigation}
Bad Data Detection (BDD) filters check the L2 norm of the measurement residuals against threshold limits.

\subsection{IDS Failure Mode}
Because the injection vector $a$ lies within the column space of the topology matrix $H$, the residual remains unchanged, allowing the attack to bypass bad data detection algorithms.

\subsection{Blockchain Mitigation}
Distributed state verification requires multiple validator nodes to sign state transitions. The topology is stored on-chain, and threshold consensus rejects transactions that violate physical law constraints.

\subsection{Mathematical Probability Model}
The probability of a successful injection is $P_{\mathrm{FDI}} = p_2$. RVR consensus mitigates this to:
\begin{equation}
P_{\mathrm{FDI,mitigated}} = (1 - \beta) P_{\mathrm{FDI}}
\end{equation}

\subsection{Simulation Parameters}
The baseline probability is $P_2 = 0.15$, with BFT mitigation factor $\beta = 0.5$.

\bigskip

% ─────────────────────────────────────────────────────────────────
\section{Communication Hijack}
\label{sec:tm:comm_hijack}
% ─────────────────────────────────────────────────────────────────

\subsection{Architecture}
This attack targets the network communication channel connecting smart meters to collectors (NAN) and utility backhauls (WAN).

\subsection{Attack Steps}
\begin{enumerate}
  \item The adversary intercepts cellular or power line communication (PLC) links.
  \item The network interface cards are compromised, or encryption keys are leaked.
  \item Transmitted packets are modified or redirected.
\end{enumerate}

\subsection{Impact}
Interception leads to data loss, leakage of customer billing details, and loss of real-time monitoring capabilities.

\subsection{Existing Mitigation}
Transport layer security protocols (TLS/IPsec) are used to encrypt and sign communication packets.

\subsection{IDS Failure Mode}
If intermediate network hops are compromised (e.g., cell tower gateways), encryption keys can be bypassed or certificates forged, rendering TLS protections ineffective.

\subsection{Blockchain Mitigation}
Data packet verification is decoupled from transport security. The ledger verifies the data origin and signature at the application layer, ensuring that manipulated packets are discarded.

\subsection{Mathematical Probability Model}
The communication hijack probability $\Pca(x)$ is:
\begin{equation}
\Pca(x) = 1 - (1 - x)^{K_2}
\end{equation}

\subsection{Simulation Parameters}
The baseline probability is $P_3 = \Pca(0.5) \approx 0.9998$, with segment count $K_2 = 1{,}271$.

\bigskip

% ─────────────────────────────────────────────────────────────────
\section{Man-in-the-Middle (MitM)}
\label{sec:tm:mitm}
% ─────────────────────────────────────────────────────────────────

\subsection{Architecture}
MitM attacks target WAN links connecting MDMS instances to utility operators, placing an interception node in the path.

\subsection{Attack Steps}
\begin{enumerate}
  \item The adversary spoofs DNS records or routes IP traffic through an intermediary node.
  \item The intermediary establishes independent connections with both endpoints.
  \item The attacker intercepts, alters, and forwards data packets.
\end{enumerate}

\subsection{Impact}
Both endpoints believe they are communicating securely, but the attacker has complete read and write access to the data stream.

\subsection{Existing Mitigation}
Public Key Infrastructure (PKI) and certificate pinning are used to verify endpoints.

\subsection{IDS Failure Mode}
If a certificate authority is compromised or root certificates are modified on client nodes, the MitM node can issue valid certificates, bypassing checks.

\subsection{Blockchain Mitigation}
Multi-signature verification requires signatures from several independent validators. A single forged certificate is insufficient to modify the ledger state.

\subsection{Mathematical Probability Model}
The probability of a successful attack is $P_4 = p_4$.

\subsection{Simulation Parameters}
The baseline probability is $P_4 = 0.05$.

\bigskip

% ─────────────────────────────────────────────────────────────────
\section{Replay Attack}
\label{sec:tm:replay}
% ─────────────────────────────────────────────────────────────────

\subsection{Architecture}
Replay attacks target local communication links (ZigBee/PLC) between smart meters and data concentrators.

\subsection{Attack Steps}
\begin{enumerate}
  \item The adversary records valid meter transmission packets containing high-demand readings.
  \item The recorded packets are retransmitted at a later time (e.g., during low-rate intervals).
  \item The concentrator processes the replayed packet as a fresh transmission.
\end{enumerate}

\subsection{Impact}
Replayed messages can trigger incorrect pricing decisions or cause billing discrepancies.

\subsection{Existing Mitigation}
Sequence numbers and time-stamp verification are checked at the receiver.

\subsection{IDS Failure Mode}
If the receiver system lacks strict clock synchronization or sequence validation limits, the replayed packet is accepted.

\subsection{Blockchain Mitigation}
The blockchain state machine prevents double-spending and replayed transactions. Each transaction is signed with a unique nonce and PoH tick height. Replayed transactions fail nonce checks and are rejected by validators.

\subsection{Mathematical Probability Model}
The probability of a successful replay is $P_5 = p_5$.

\subsection{Simulation Parameters}
The baseline probability is $P_5 = 0.15$.

\bigskip

% ─────────────────────────────────────────────────────────────────
\section{Denial of Service (DoS)}
\label{sec:tm:dos}
% ─────────────────────────────────────────────────────────────────

\subsection{Architecture}
DoS attacks target concentrators and utility collectors in NAN/WAN regions.

\subsection{Attack Steps}
\begin{enumerate}
  \item The adversary floods network interfaces with malformed packets.
  \item The processing buffers on target devices are exhausted.
  \item The target device drops incoming connection requests.
\end{enumerate}

\subsection{Impact}
Utility operators lose real-time visibility into grid metrics and control commands are delayed.

\subsection{Existing Mitigation}
Firewalls and rate-limiting limits are applied at the network edge.

\subsection{IDS Failure Mode}
Static rate limits can block legitimate meter traffic during peak events, and fail to prevent distributed resource exhaustion.

\subsection{Blockchain Mitigation}
Distributed ledger systems eliminate single-point-of-failure endpoints. If a validator node is targeted by a DoS attack, the consensus group continues to operate with the remaining nodes.

\subsection{Mathematical Probability Model}
The probability of a successful DoS attack is $P_6 = p_6$.

\subsection{Simulation Parameters}
The baseline probability is $P_6 = 0.20$.

\bigskip

% ─────────────────────────────────────────────────────────────────
\section{Distributed Denial of Service (DDoS)}
\label{sec:tm:ddos}
% ─────────────────────────────────────────────────────────────────

\subsection{Architecture}
DDoS attacks target centralized servers, such as the utility MDMS, using distributed botnets.

\subsection{Attack Steps}
\begin{enumerate}
  \item The adversary compromises a large number of IoT endpoints.
  \item The compromised nodes send coordinated traffic to the central target.
  \item The central network link is overwhelmed by the high volume of incoming data.
\end{enumerate}

\subsection{Impact}
Centralized databases become unavailable, disrupting billing operations, load balancing, and market clearing.

\subsection{Existing Mitigation}
DDoS scrubbing centers and traffic filtering are used to clean incoming data.

\subsection{IDS Failure Mode}
High-volume attacks can overwhelm scrubbing centers, and centralized systems remain vulnerable to complete resource exhaustion.

\subsection{Blockchain Mitigation}
Ledger state replication distributes the communication load across multiple consensus nodes. High-throughput protocols (e.g., Tower BFT) use network topology trees to handle transaction routing, reducing the impact of distributed flooding.

\subsection{Mathematical Probability Model}
The probability of a successful DDoS attack is $P_7 = p_7$.

\subsection{Simulation Parameters}
The baseline probability is $P_7 = 0.35$.

\bigskip

% ─────────────────────────────────────────────────────────────────
\section{Key Compromise}
\label{sec:tm:key_compromise}
% ─────────────────────────────────────────────────────────────────

\subsection{Architecture}
This attack targets the private cryptographic keys stored on smart meters and validator nodes.

\subsection{Attack Steps}
\begin{enumerate}
  \item The adversary exploits memory disclosure vulnerabilities or uses physical side-channel attacks.
  \item The private key material is extracted from the device memory.
  \item The attacker uses the key to sign forged transactions.
\end{enumerate}

\subsection{Impact}
The attacker gains full control over the compromised node's identity, allowing them to inject validly signed false data.

\subsection{Existing Mitigation}
Hardware security modules (HSM) and secure enclaves are used to protect key material.

\subsection{IDS Failure Mode}
Once a private key is compromised, any data signed with that key appears legitimate to the IDS, bypassing signature checks.

\subsection{Blockchain Mitigation}
Threshold cryptography (e.g., Shamir Secret Sharing) is used to divide the key material. Compromising a single validator share does not expose the master signing key, which requires a threshold number of shares to reconstruct.

\subsection{Mathematical Probability Model}
The key compromise probability is:
\begin{equation}
P_{\mathrm{KC,mitigated}} = \eta \cdot P_8
\end{equation}

\subsection{Simulation Parameters}
The baseline probability is $P_8 = 0.01$, with threshold mitigation factor $\eta = 0.01$ (RVR) or $\eta = 1.0$ (PBFT).

\bigskip

% ─────────────────────────────────────────────────────────────────
\section{SCADA Breach}
\label{sec:tm:scada_breach}
% ─────────────────────────────────────────────────────────────────

\subsection{Architecture}
This attack targets the database servers inside the utility control center.

\subsection{Attack Steps}
\begin{enumerate}
  \item The adversary penetrates the corporate network using stolen credentials.
  \item The attacker moves laterally into the SCADA network zone.
  \item Direct SQL injection or database command execution is used to modify grid configuration files.
\end{enumerate}

\subsection{Impact}
The utility database is corrupted, potentially leading to incorrect control commands being sent to field actuators.

\subsection{Existing Mitigation}
Firewalls and database access control lists (ACL) are used to restrict access.

\subsection{IDS Failure Mode}
If administrative credentials are compromised, the security logs and historical records can be modified by the attacker to cover their tracks.

\subsection{Blockchain Mitigation}
The immutable design of the blockchain ledger prevents history modification. Since the data is cryptographically chained, any retrospective modification changes the hash history, which is detected and rejected by other validators.

\subsection{Mathematical Probability Model}
The probability of a SCADA breach is $P_9 = p_9$.

\subsection{Simulation Parameters}
The baseline probability is $P_9 = 0.01$.

\bigskip

% ─────────────────────────────────────────────────────────────────
\section{Receiver Override}
\label{sec:tm:receiver_override}
% ─────────────────────────────────────────────────────────────────

\subsection{Architecture}
This attack targets client software interfaces and automated billing systems.

\subsection{Attack Steps}
\begin{enumerate}
  \item The adversary gains unauthorized access to utility billing consoles.
  \item Compromised instructions are sent to smart meters to execute shut-off commands.
  \item System logs are cleared to delay response.
\end{enumerate}

\subsection{Impact}
Widespread power shutdowns can be triggered remotely, disrupting customer service.

\subsection{Existing Mitigation}
Role-Based Access Control (RBAC) is used to manage system permissions.

\subsection{IDS Failure Mode}
If the administrator account is compromised, the attacker has full authorization to issue override commands, bypassing standard security filters.

\subsection{Blockchain Mitigation}
Critical control commands require multi-signature verification from multiple utility operators. An override command will not be processed unless it receives the threshold number of signatures.

\subsection{Mathematical Probability Model}
The probability of a receiver override is $P_{10} = p_{10}$.

\subsection{Simulation Parameters}
The baseline probability is $P_{10} = 0.01$.

\bigskip

% ─────────────────────────────────────────────────────────────────
\section{Sybil Attack}
\label{sec:tm:sybil}
% ─────────────────────────────────────────────────────────────────

\subsection{Architecture}
Sybil attacks target peer-to-peer communication layers and validator election systems in distributed networks.

\subsection{Attack Steps}
\begin{enumerate}
  \item The adversary creates multiple virtual nodes on a single physical host.
  \item The virtual nodes register with the network using fabricated identities.
  \item The Sybil nodes act in coordination to influence validator elections or transaction routing.
\end{enumerate}

\subsection{Impact}
The adversary can gain a majority in consensus quorums, allowing them to control block generation and double-spend transactions.

\subsection{Existing Mitigation}
IP address filtering and identity registration are checked.

\subsection{IDS Failure Mode}
In large AMI deployments, IP filtering can block legitimate devices using dynamic addresses, and fail to prevent distributed virtual nodes from registering.

\subsection{Blockchain Mitigation}
Permissioned validator registries verify node identities before admission. In RVR consensus, leader selection is weighted by credit scoring, which makes it harder for newly created Sybil nodes without transaction history to influence the network.

\subsection{Mathematical Probability Model}
The probability of a successful Sybil attack is:
\begin{equation}
P_{\mathrm{Sybil,mitigated}} = \omega_{\mathrm{credit}} \cdot P_{11}
\end{equation}

\subsection{Simulation Parameters}
The baseline probability is $P_{11} = 0.08$, with credit weight factor $\omega_{\mathrm{credit}} = 0.001$ (RVR) or $\omega_{\mathrm{credit}} = 1.0$ (PBFT).

\bigskip

% ─────────────────────────────────────────────────────────────────
\section{Byzantine Validator Compromise}
\label{sec:tm:byz_compromise}
% ─────────────────────────────────────────────────────────────────

\subsection{Architecture}
This attack targets the active consensus validators in a distributed ledger network.

\subsection{Attack Steps}
\begin{enumerate}
  \item The adversary compromises a subset of validator nodes.
  \item The compromised nodes send conflicting messages (double-voting) during consensus rounds.
  \item The attacker attempts to stall the consensus process or commit invalid blocks.
\end{enumerate}

\subsection{Impact}
If the number of compromised nodes exceeds the threshold limit, the ledger safety is breached, potentially leading to double-spending or network forks.

\subsection{Existing Mitigation}
Centralized node monitoring and log verification are used to detect anomalies.

\subsection{IDS Failure Mode}
Centralized monitors cannot prevent a quorum of compromised nodes from committing a block once consensus is reached.

\subsection{Blockchain Mitigation}
Consensus protocols are designed to tolerate up to $f = \lfloor (n-1)/3 \rfloor$ faulty nodes. By using validator rotation and VRF selection, the probability that the adversary compromises a consensus quorum is minimized.

\subsection{Mathematical Probability Model}
The probability of consensus failure is modeled by the binomial tail:
\begin{equation}
\Pbyz(n, f, p_c) = \sum_{k=f+1}^{n} \binom{n}{k} p_c^k (1 - p_c)^{n - k}
\end{equation}

\subsection{Simulation Parameters}
The baseline probability is $P_{12} = 0.05$, with node count $n = 51$ and threshold limit $f = 16$.

\bigskip

% ─────────────────────────────────────────────────────────────────
\section{Chapter Summary}
\label{sec:tm:summary}
% ─────────────────────────────────────────────────────────────────

We summarize the key details of the twelve attack vectors in Table~\ref{tab:attack_comparison_matrix}.

\begin{table}[htbp]
\centering
\small
\renewcommand{\arraystretch}{1.3}
\begin{tabularx}{\textwidth}{lXll}
\toprule
\textbf{Attack Vector} & \textbf{Core Target} & \textbf{IDS Failure Mode} & \textbf{Blockchain Mitigation} \\
\midrule
Sensor Compromise & Meter firmware & Slow reading drift & Immutable record signing \\
FDI & State estimator & Topology bypass & Consensus rule checks \\
Comm Hijack & Network segments & Key decryption & Origin authentication \\
MitM & WAN links & Certificate forgery & Multi-signature approvals \\
Replay & Local message links & Lacks sequence checks & Nonce & time-stamp limits \\
DoS & Network interfaces & Central resource exhaust & Node redundancy \\
DDoS & Utility MDMS & scrubbing overload & Distributed replicas \\
Key Compromise & Device memory & Signed fake data is accepted & Shamir Secret Sharing \\
SCADA Breach & centralized database & Administrative override & Cryptographic chaining \\
Receiver Override & utility billing console & Master account takeover & Multi-operator approval \\
Sybil & P2P registry & IP verification limit & Credit-weighted selection \\
Byz Validator & consensus nodes & Cannot stop commit quorum & Byzantine threshold ($f < n/3$) \\
\bottomrule
\end{tabularx}
\caption{Twelve-attack comparative threat evaluation matrix detailing targets, IDS limits, and BFT mitigations.}
\label{tab:attack_comparison_matrix}
\end{table}

% --- END INPUT: chapters/ch03_threat_model ---

\cleardoublepage


% --- BEGIN INPUT: chapters/ch04_security_framework ---
% ══════════════════════════════════════════════════════════════════
% CHAPTER 7: SECURITY MATHEMATICS
% ══════════════════════════════════════════════════════════════════
\chapter{Security Mathematics}
\label{ch:security_mathematics}

This chapter establishes the analytical mathematical security framework that governs the comparison between reactive detection and proactive consensus systems in Advanced Metering Infrastructure (AMI). Every derivation is presented from first principles, proceeding systematically from general assumptions to theorems, formal proofs, and clear physical interpretations. All mathematical assertions are purely analytical \analytical.

\section{Probability Theory}
\label{sec:sm:probability}

We define our framework on a standard Kolmogorov probability space $(\Omega, \mathcal{F}, P)$, where $\Omega$ is the sample space of system vulnerability states, $\mathcal{F}$ is the $\sigma$-algebra of events, and $P$ is the probability measure mapping $\mathcal{F} \to [0, 1]$.

\begin{definition}[Compromise State]
Let $Y_j \in \{0, 1\}$ be a Bernoulli random variable representing the compromise state of the $j$-th cyber-physical asset, where:
\begin{equation}
Y_j = \begin{cases}
1, & \text{if asset } j \text{ is successfully compromised} \\
0, & \text{otherwise}
\end{cases}
\end{equation}
The marginal probability of compromise is $P(Y_j = 1) = P(\mathcal{C}_j) = p_j$, where $\mathcal{C}_j$ is the compromise event.
\end{definition}

\subsection{Conditional and Joint Events}
For any two attack events $\mathcal{C}_i$ and $\mathcal{C}_j$, their conditional relation models coordinated intrusion:
\begin{equation}
P(\mathcal{C}_i \mid \mathcal{C}_j) = \frac{P(\mathcal{C}_i \cap \mathcal{C}_j)}{P(\mathcal{C}_j)}, \quad P(\mathcal{C}_j) > 0
\end{equation}
When the adversary coordinates attacks, they create dependency. If they act independently, the joint probability simplifies to:
\begin{equation}
P(\mathcal{C}_i \cap \mathcal{C}_j) = P(\mathcal{C}_i) \cdot P(\mathcal{C}_j) = p_i p_j
\end{equation}

\subsection{Bayesian Updates under Detection}
When an intrusion detection alarm is triggered (event $\mathcal{A}$), the posterior probability that an asset is actually compromised is updated using Bayes' rule:
\begin{equation}
P(\mathcal{C}_j \mid \mathcal{A}) = \frac{P(\mathcal{A} \mid \mathcal{C}_j) P(\mathcal{C}_j)}{P(\mathcal{A})} = \frac{D_j p_j}{D_j p_j + F_j (1 - p_j)}
\end{equation}
where $D_j = P(\mathcal{A} \mid \mathcal{C}_j)$ is the detection rate (true positive rate) and $F_j = P(\mathcal{A} \mid \overline{\mathcal{C}_j})$ is the false alarm rate (false positive rate). Under high false alarm rates ($F_j \gg 0$), the operational confidence $P(\mathcal{C}_j \mid \mathcal{A})$ degrades, leading to operator fatigue and reactive delay.

\section{Reliability Theory}
\label{sec:sm:reliability}

In security engineering, we model the system defense using reliability structures. An AMI security system under a multi-vector threat landscape functions as a parallel-failure (or series-defense) reliability system.

\begin{definition}[System Survival]
The system survives if and only if no attack event $\mathcal{C}_j$ succeeds. The system compromise event $\mathcal{S}$ is:
\begin{equation}
\mathcal{S} = \bigcup_{j=1}^{m} \mathcal{C}_j
\end{equation}
\end{definition}

\begin{assumption}[Vulnerability Independence]
The compromise events $\{\mathcal{C}_j\}_{j=1}^m$ across different vectors are mutually independent.
\end{assumption}

\begin{lemma}[Reliability Product]
Under the independent parallel failure model, the joint compromise probability $P(\mathcal{S})$ is:
\begin{equation}
P(\mathcal{S}) = 1 - \prod_{j=1}^{m} (1 - p_j)
\end{equation}
\end{lemma}

\begin{proof}
Let $\overline{\mathcal{S}}$ be the survival event. Then $\overline{\mathcal{S}} = \bigcap_{j=1}^m \overline{\mathcal{C}_j}$. By independence, $P(\overline{\mathcal{S}}) = \prod_{j=1}^m P(\overline{\mathcal{C}_j}) = \prod_{j=1}^m (1 - p_j)$. Since $P(\mathcal{S}) = 1 - P(\overline{\mathcal{S}})$, the result follows.
\end{proof}

\section{Random Processes}
\label{sec:sm:processes}

To represent the dynamic arrival of cyber threats in real-time grid operations, we model the threat landscape as a stochastic arrival process.

\begin{assumption}[Homogeneous Poisson Arrivals]
Cyber-physical attacks arrive at the AMI system according to a homogeneous Poisson process $\{N(t), t \ge 0\}$ with rate parameter $\lambda > 0$ attacks per second.
\end{assumption}

\begin{definition}[Poisson Arrival Distribution]
The probability of receiving exactly $k$ attacks in a time window of duration $\tau$ is:
\begin{equation}
P(N(t+\tau) - N(t) = k) = \frac{(\lambda \tau)^k e^{-\lambda \tau}}{k!}, \quad k = 0, 1, 2, \ldots
\end{equation}
\end{definition}
The inter-arrival times $T_k$ between successive attacks are independent and identically distributed exponential random variables:
\begin{equation}
f_{T}(t) = \lambda e^{-\lambda t}, \quad t \ge 0
\end{equation}

\section{Static Security Model}
\label{sec:sm:static_model}

This section reproduces the baseline security model of Sheikh et al. \cite{sheikh2020}. The model analyzes static compromise probabilities across four main sectors of an AMI system:
\begin{enumerate}
  \item \textbf{Sensor Level ($\Psa$):} Vulnerability in measurement units, modeled under pairwise coordinated compromise.
  \item \textbf{Communication Level ($\Pca$):} Hijacking of intermediate networking hops.
  \item \textbf{SCADA Level ($\Pscada$):} Direct compromise of the centralized control database.
  \item \textbf{Receiver Level ($\Pr$):} Compromise of utility terminals.
\end{enumerate}

\begin{assumption}[IEEE 33-Bus Configuration]
The reference system contains $N_{\mathrm{sen}} = 3{,}854$ sensors. The intermediate communication is divided into $K_2 = 1{,}271$ independent network segments. SCADA and receiver systems maintain fixed baseline risk levels: $\Pscada = 0.01$ and $\Pr = 0.01$ \litderived.
\end{assumption}

\begin{lemma}[Pairwise Sensor Probability]
The probability $\Psa(x)$ of a successful pairwise sensor compromise under system vulnerability index $x \in [0, 1]$ is:
\begin{equation}
\Psa(x) = x^2 (1 - x)^{N_{\mathrm{sen}} - 2}
\end{equation}
\end{lemma}

\begin{proof}
Let $K_1 = \binom{N_{\mathrm{sen}}}{2}$ be the total sensor pairs. The probability that exactly one pair of sensors is compromised while all others remain secure is:
\begin{equation}
P(\text{Pair}) = K_1 x^2 (1 - x)^{N_{\mathrm{sen}} - 2}
\end{equation}
Distributing the joint event across the pair combinations yields the normalized marginal probability:
\begin{equation}
\Psa(x) = \frac{K_1 x^2 (1 - x)^{N_{\mathrm{sen}} - 2}}{K_1} = x^2 (1 - x)^{N_{\mathrm{sen}} - 2}
\end{equation}
\end{proof}

The communication compromise probability is:
\begin{equation}
\Pca(x) = 1 - (1 - x)^{K_2}
\end{equation}
The joint static baseline compromise probability $\PTA(x)$ is defined as:
\begin{equation}
\PTA(x) = 1 - \bigl(1 - \Psa(x)\bigr)\bigl(1 - \Pca(x)\bigr)(1 - \Pscada)(1 - \Pr)
\end{equation}
For $x=0.5$, this results in $\PTA(0.5) \approx 0.9802$, reproducing the baseline validation point.

\section{Twelve-Attack Model}
\label{sec:sm:twelve_attack}

We expand the four aggregated baseline sectors into twelve distinct cyber-physical attack categories mapping the complete AMI vulnerability surface.

\begin{definition}[Twelve-Attack Parallel System]
Let the system be exposed to twelve independent attack probabilities $\{P_j\}_{j=1}^{12}$, corresponding to: sensor compromise, FDI, communication hijack, MitM, replay, DoS, DDoS, key compromise, SCADA breach, receiver override, Sybil, and Byzantine validator compromise.
\end{definition}

Under a reactive centralized IDS, the joint static compromise probability is:
\begin{equation}
\Pids = 1 - \prod_{j=1}^{12} (1 - P_j)
\end{equation}
Using standard parameter calibrations from Table~\ref{tab:contribution_mapping}, the joint baseline static probability evaluates to:
\begin{equation}
\Pids \approx 0.99998
\end{equation}
This indicates that under a parallel threat vector model, reactive detection systems suffer from structural failure because the failure of any single detection filter compromises the system's overall integrity.

\section{Byzantine Model}
\label{sec:sm:byzantine_model}

To evaluate the proactive consensus paradigm, we replace centralized database storage with a distributed validator set of size $n$, yielding Byzantine fault tolerance threshold bounds.

\begin{assumption}[Validator Compromise Independence]
Each of the $n$ validator nodes is compromised independently with probability $p_c \in [0, 1]$.
\end{assumption}

\begin{theorem}[Byzantine Failure Boundary]
The probability that a Byzantine Fault Tolerant consensus network with $n$ nodes and threshold $f = \lfloor (n-1)/3 \rfloor$ fails is:
\begin{equation}
\Pbyz(n, f, p_c) = \sum_{k=f+1}^{n} \binom{n}{k} p_c^k (1 - p_c)^{n - k}
\end{equation}
\end{theorem}

\begin{proof}
Let $K$ be a binomial random variable representing the number of compromised validators, where $K \sim \mathrm{Binomial}(n, p_c)$. The consensus network maintains safety and liveness if and only if $K \le f$. The probability of consensus compromise is the tail of the binomial distribution:
\begin{equation}
P(K > f) = \sum_{k=f+1}^{n} P(K=k) = \sum_{k=f+1}^{n} \binom{n}{k} p_c^k (1 - p_c)^{n - k}
\end{equation}
which matches the theorem.
\end{proof}

Under the reference calibration of $n=51$, $f=16$, and $p_c=0.05$, the compromise probability evaluates to:
\begin{equation}
\Pbyz(51, 16, 0.05) \approx 2.14 \times 10^{-10}
\end{equation}
This is an improvement of over 170 orders of magnitude compared to reactive systems under the static model.

\section{Temporal Model}
\label{sec:sm:temporal_model}

To incorporate dynamic threat behavior, we define the exposure window for both paradigms.

\begin{definition}[Vulnerability Exposure Window]
For a security system with response latency $\tau$ exposed to Poisson attack arrivals with rate $\lambda$, the probability $\Wvuln(\lambda, \tau)$ that at least one attack arrives during the response window is:
\begin{equation}
\Wvuln(\lambda, \tau) = 1 - e^{-\lambda \tau}
\end{equation}
\end{definition}

This represents the time window during which the system is exposed before mitigation is finalized.
\begin{itemize}
  \item \textbf{Reactive IDS Window:} $\tau_{\mathrm{IDS}} = 30$~s $\implies \Wvuln = 1 - e^{-30\lambda}$
  \item \textbf{Proactive RVR Window:} $\tau_{\mathrm{RVR}} = 0.2$~s $\implies \Wvuln = 1 - e^{-0.2\lambda}$
\end{itemize}

\section{Combined Index}
\label{sec:sm:combined_index}

We define the Static-Temporal Security Index (STSI) to combine structural vulnerability and dynamic exposure into a single metric.

\begin{definition}[Static-Temporal Security Index]
The STSI of a system with static compromise probability $P_{\mathrm{static}}$ and response latency $\tau$ under threat rate $\lambda$ is:
\begin{equation}
\mathrm{STSI}(\lambda, \tau, P_{\mathrm{static}}) = P_{\mathrm{static}} \times \bigl(1 - e^{-\lambda \tau}\bigr)
\end{equation}
\end{definition}

The STSI provides a clear comparison of security paradigms. For RVR consensus:
\begin{equation}
\mathrm{STSI}_{\mathrm{RVR}} = \Pbyz(51, 16, 0.05) \times \bigl(1 - e^{-0.2\lambda}\bigr) \approx 2.14 \times 10^{-10} \times \bigl(1 - e^{-0.2\lambda}\bigr)
\end{equation}
This formulation demonstrates that distributed consensus reduces overall risk exposure by combining Byzantine threshold security with sub-second response times.

% Consensus Mathematics Section
% ══════════════════════════════════════════════════════════════════
% CHAPTER 8: CONSENSUS MATHEMATICS
% ══════════════════════════════════════════════════════════════════



This chapter provides the mathematical and algorithmic formulation of the Byzantine Fault Tolerant (BFT) consensus protocols analyzed in this thesis. All protocol descriptions, complexity derivations, and safety/liveness proofs are analytical \analytical\ and follow a standardized layout to facilitate comparison.

The evaluated consensus protocols span four evolutionary generations:
\begin{enumerate}
  \item \textbf{G1: Classical BFT} (Classic PBFT, Oral Message OM(m))
  \item \textbf{G2: Committee-Delegated BFT} (IBFT 2.0, QBFT)
  \item \textbf{G3: Hierarchical BFT} (CE-PBFT, G-PBFT, SV-PBFT)
  \item \textbf{G4: Sub-Second BFT} (Tower BFT, RVR)
\end{enumerate}

\bigskip

% ─────────────────────────────────────────────────────────────────
\chapter{Consensus Mathematics}
\label{ch:consensus_mathematics}

This chapter details the formal mathematical algorithms and models underlying the consensus mechanisms evaluated in this research.

\section{Classic PBFT}
\label{sec:cm:pbft}
% ─────────────────────────────────────────────────────────────────

\subsection{Overview}
Practical Byzantine Fault Tolerance (PBFT) was introduced by Castro and Liskov \cite{castro1999} to provide state machine replication in asynchronous environments containing Byzantine failures.

\subsection{Assumptions}
\begin{enumerate}
  \item The network operates under partial synchrony, where message delays are bounded by an unknown $\Delta$ after a Global Stabilization Time (GST).
  \item The validator set size is $n$, with a maximum of $f = \lfloor (n-1)/3 \rfloor$ Byzantine nodes.
  \item Cryptographic primitives (digital signatures and hash functions) are secure against polynomial-time adversaries.
\end{enumerate}

\subsection{Algorithm}
\begin{algorithm}[H]
\caption{Classic PBFT Consensus Round}
\label{alg:pbft}
\begin{algorithmic}[1]
\STATE \textbf{Input:} Validator set $\mathcal{V}$ of size $n$, current view $v$, transaction batch $B$
\STATE \textbf{Output:} Committed block $Block$ appended to the ledger
\STATE \textbf{Phase 1: Pre-Prepare}
\IF{validator is primary $p = v \bmod n$}
    \STATE Broadcast $\langle\text{PRE-PREPARE}, v, n_{\mathrm{seq}}, d\rangle_{\sigma_p}$ and $B$ to $\mathcal{V}$
\ENDIF
\STATE \textbf{Phase 2: Prepare}
\IF{validator $i$ receives valid PRE-PREPARE from $p$}
    \STATE Verify hash $d = \mathrm{Hash}(B)$ and view $v$
    \STATE Broadcast $\langle\text{PREPARE}, v, n_{\mathrm{seq}}, d, i\rangle_{\sigma_i}$ to $\mathcal{V}$
\ENDIF
\STATE \textbf{Phase 3: Commit}
\IF{validator $i$ collects a prepared certificate (1 PRE-PREPARE and $2f$ PREPARE messages matching $v, n_{\mathrm{seq}}, d$)}
    \STATE Broadcast $\langle\text{COMMIT}, v, n_{\mathrm{seq}}, d, i\rangle_{\sigma_i}$ to $\mathcal{V}$
\ENDIF
\STATE \textbf{Phase 4: Append}
\IF{validator $i$ collects a commit certificate ($2f+1$ COMMIT messages matching $v, n_{\mathrm{seq}}, d$)}
    \STATE Execute transactions in $B$ and append block to local ledger
\ENDIF
\end{algorithmic}
\end{algorithm}

\subsection{Message Complexity}
In the three-phase consensus process (Pre-prepare, Prepare, Commit), each of the $n$ nodes broadcasts its Prepare and Commit messages to all other $n-1$ nodes.
\begin{equation}
M_{\mathrm{PBFT}} = 1 \times (n-1) \quad (\text{Pre-Prepare}) + n \times (n-1) \quad (\text{Prepare}) + n \times (n-1) \quad (\text{Commit})
\end{equation}
This yields a quadratic message complexity:
\begin{equation}
M_{\mathrm{PBFT}} = 2n^2 - n - 1 \in \mathcal{O}(n^2)
\end{equation}

\subsection{Latency}
Let $\tau_{\mathrm{net}}$ represent the average network propagation delay. The three communication steps require:
\begin{equation}
\tau_{\mathrm{PBFT}} = 3 \tau_{\mathrm{net}} + \tau_{\mathrm{proc}}
\end{equation}
Under our baseline calibrations (with queuing delays and cryptographic verification), this results in an empirical round latency of $\tau_{\mathrm{PBFT}} \approx 7{,}650.0$~ms.

\subsection{Safety}
Safety guarantees that honest nodes do not commit conflicting blocks at the same sequence number.
\begin{theorem}[PBFT Safety]
If two honest nodes commit blocks $B$ and $B'$ at the same sequence number $n_{\mathrm{seq}}$, then $B = B'$.
\end{theorem}
\begin{proof}
For a node to commit a block, it must collect a commit certificate containing $2f+1$ COMMIT messages. Let $\mathcal{Q}_1$ and $\mathcal{Q}_2$ represent the quorum sets of validators that sent COMMIT messages for $B$ and $B'$ respectively. Since $|\mathcal{Q}_1| \ge 2f+1$ and $|\mathcal{Q}_2| \ge 2f+1$ from a total set of $n = 3f+1$ validators:
\begin{equation}
|\mathcal{Q}_1 \cap \mathcal{Q}_2| \ge (2f+1) + (2f+1) - (3f+1) = f+1
\end{equation}
Since there are at most $f$ Byzantine nodes in the network, the intersection must contain at least one honest validator. An honest validator only commits a single block per sequence number in a given view. Thus, the certificates must match, proving $B = B'$.
\end{proof}

\subsection{Liveness}
Liveness guarantees that the system continues to make progress and commit transactions. Under PBFT, liveness is maintained by the View Change protocol. If the primary node fails or is compromised, a timer expires, and validators broadcast view-change messages. A new view is established when $2f+1$ nodes agree on the transition. Because the network is partially synchronous, liveness is guaranteed after GST.

\subsection{Fault Tolerance}
The maximum number of Byzantine failures the system can tolerate is:
\begin{equation}
f < \frac{n}{3}
\end{equation}
This threshold is optimal for non-synchronous consensus networks.

\subsection{AMI Suitability}
PBFT provides immediate finality, ensuring that smart meter transactions cannot be forks or reversed. However, its $\mathcal{O}(n^2)$ message complexity limits its scalability to validator networks of size $n \le 100$, making it unsuitable for direct smart meter deployment.

\subsection{Limitations}
\begin{enumerate}
  \item Quadratic message overhead prevents scaling to large validator sets.
  \item High view-change overhead under adversarial conditions.
  \item Direct vulnerability to single primary node DoS attacks.
\end{enumerate}

\bigskip

% ─────────────────────────────────────────────────────────────────
\section{Tower BFT}
\label{sec:cm:tower}
% ─────────────────────────────────────────────────────────────────

\subsection{Overview}
Tower BFT is a sub-second consensus protocol introduced by Yakovenko \cite{yakovenko2018} for high-performance networks, utilizing a verifiable clock known as Proof-of-History (PoH) to order transactions.

\subsection{Assumptions}
\begin{enumerate}
  \item The system maintains a synchronized PoH ledger serving as a global clock.
  \item The validator set size is $n$, with Byzantine fraction $f \le \lfloor (n-1)/3 \rfloor$.
  \item Cryptographic verification is accelerated via GPU parallel processing.
\end{enumerate}

\subsection{Algorithm}
\begin{algorithm}[H]
\caption{Tower BFT Consensus Vote}
\label{alg:tower}
\begin{algorithmic}[1]
\STATE \textbf{Input:} Validator set $\mathcal{V}$ of size $n$, local block queue $\mathcal{Q}$, PoH tick $t$
\STATE \textbf{Output:} Vote transaction submitted to the PoH log
\STATE \textbf{Phase 1: Propose}
\STATE Primary node packaging block $B$ hashes it into the active PoH chain at tick $t$
\STATE Broadcast $B$ to $\mathcal{V}$
\STATE \textbf{Phase 2: Vote Selection}
\IF{validator $i$ receives $B$ and verifies PoH sequence}
    \STATE Compute vote timeout $d = 2^k$ slots (where $k$ is the vote stack depth)
    \STATE Sign and submit $\langle\text{VOTE}, B, i, t\rangle_{\sigma_i}$ as a transaction
\ENDIF
\STATE \textbf{Phase 3: Confirmation}
\IF{block $B$ collects votes from a supermajority ($2f+1$ validators)}
    \STATE Lock the block on the Tower; rollback depth increases exponentially
\ENDIF
\end{algorithmic}
\end{algorithm}

\subsection{Message Complexity}
By utilizing the PoH clock for scheduling and ordering, validators submit votes directly to the ledger leader rather than broadcasting to all nodes. The message complexity is linear:
\begin{equation}
M_{\mathrm{Tower}} = n \in \mathcal{O}(n)
\end{equation}


\subsection{Properties}
\begin{itemize}
    \item \textbf{Safety:} Maintains safety up to  = \lfloor (n-1)/3 \rfloor$ malicious nodes, relying on threshold signatures to aggregate votes and verifiable random functions to randomly select a committee.
    \item \textbf{Liveness:} Guaranteed by rotating the primary through a verifiable random function mechanism, reducing primary censorship.
    \item \textbf{Termination:} Transactions are finalized deterministically once the threshold signature is assembled for a block.
\end{itemize}

\subsection{Latency}
PoH allows pipelined block generation, reducing the consensus delay to:
\begin{equation}
\tau_{\mathrm{Tower}} = \tau_{\mathrm{slot}} \approx 242.9 \text{ ms}
\end{equation}

\subsection{Safety}
Safety is enforced by lockouts. When a validator votes for a block on the tower, it commits to a lockout time $2^k$ slots, during which it cannot vote for any block that does not descend from the voted block. If a fork occurs, the validator must wait for the lockout to expire before voting on the competing branch. This prevents split-brain forks, ensuring safety.

\subsection{Liveness}
Liveness is guaranteed because the exponential lockouts prevent validators from getting stuck in deadlocks on minor forks. Once a block gathers a supermajority of locks, the rollback time exceeds the lifespan of the network, finalizing the state.

\subsection{Fault Tolerance}
Tower BFT maintains standard Byzantine resistance:
\begin{equation}
f = \left\lfloor \frac{n-1}{3} \right\rfloor
\end{equation}

\subsection{AMI Suitability}
The sub-second latency ($\approx 240$~ms) and linear scaling make Tower BFT highly suitable for large-scale utility MDMS and HES systems handling high-frequency smart meter readings.

\subsection{Limitations}
\begin{enumerate}
  \item High dependency on the stability of the PoH leader generator.
  \item Leader-rotation overhead can cause transaction drops.
  \item Requires substantial hardware resources (GPUs) for signature validation.
\end{enumerate}

\bigskip

% ─────────────────────────────────────────────────────────────────
\section{RVR Consensus}
\label{sec:cm:rvr}
% ─────────────────────────────────────────────────────────────────

\subsection{Overview}
Random Verifiable Rotation (RVR) consensus, proposed by Wang et al. \cite{wang2025}, is a weighted Byzantine-tolerant consensus algorithm designed for smart grid trading. It uses credit scoring and a Verifiable Random Function (VRF) to select consensus groups.

\subsection{Assumptions}
\begin{enumerate}
  \item Each validator has a credit score $C_i \in [0, 100]$ updated based on past behavior.
  \item A verifiable random seed is generated via VRF.
  \item The validator set size is $n$, with a maximum of $f = \lfloor (n-1)/3 \rfloor$ Byzantine nodes.
\end{enumerate}

\subsection{Algorithm}
\begin{algorithm}[H]
\caption{RVR Leader Selection and Consensus}
\label{alg:rvr}
\begin{algorithmic}[1]
\STATE \textbf{Input:} Validator set $\mathcal{V}$ of size $n$, credit vector $\vec{C}$, VRF seed $S$
\STATE \textbf{Output:} Selected consensus group $\mathcal{G}_{\mathrm{c}}$ and committed transaction
\STATE \textbf{Phase 1: VRF Selection}
\STATE Each node computes $v_i = \mathrm{VRF}_{\mathrm{sk}_i}(S)$ and proof $\pi_i$
\STATE Select consensus leader $L$ based on highest weighted value: $w_i = v_i \times C_i$
\STATE \textbf{Phase 2: Group Election}
\STATE Leader $L$ selects a subset of validators $\mathcal{G}_{\mathrm{c}}$ of size $m \ll n$ based on credit rank
\STATE \textbf{Phase 3: Consensus Execution}
\STATE Execute a simplified single-phase PBFT round within group $\mathcal{G}_{\mathrm{c}}$
\STATE Broadcast committed state to the main validator set $\mathcal{V}$
\end{algorithmic}
\end{algorithm}

\subsection{Message Complexity}
By delegating the active consensus phase to a small group $\mathcal{G}_{\mathrm{c}}$ of size $m$, the overall message complexity is reduced to:
\begin{equation}
M_{\mathrm{RVR}} = m^2 + n \in \mathcal{O}(m^2 + n)
\end{equation}
Because $m \ll n$, this scales linearly with the total number of nodes.

\subsection{Latency}
RVR minimizes round complexity to:
\begin{equation}
\tau_{\mathrm{RVR}} \approx 200.0 \text{ ms}
\end{equation}

\subsection{Safety}
Safety is guaranteed within the consensus group $\mathcal{G}_{\mathrm{c}}$ through threshold signatures. The credit ranking ensures that nodes with a history of Byzantine behavior are excluded from the consensus group, lowering the probability of selecting $f+1$ faulty validators.

\subsection{Liveness}
Liveness is maintained by rotating the VRF seed and consensus group every block round. If a selected consensus group fails to reach agreement, the current slot expires, a new VRF seed is generated, and a new group is selected. This prevents permanent deadlocks.

\subsection{Fault Tolerance}
RVR tolerates standard Byzantine conditions $f < n/3$, but incorporates credit-based weights that lower the probability of Byzantine failure.

\subsection{AMI Suitability}
RVR's combination of credit-based Sybil resistance and sub-second latency makes it highly suitable for peer-to-peer energy trading and secure data collection in AMI systems.

\subsection{Limitations}
\begin{enumerate}
  \item Credit calculation introduces storage overhead.
  \item Vulnerability to credit-boosting attacks where nodes act honestly to gain high rank before launching a coordinated attack.
\end{enumerate}

\bigskip

% ─────────────────────────────────────────────────────────────────
\section{Summary and Comparative Taxonomy}
\label{sec:cm:summary}
% ─────────────────────────────────────────────────────────────────

We summarize the mathematical and performance parameters of the nine BFT consensus protocols in Table~\ref{tab:consensus_parameters_full}.

\begin{table}[htbp]
\centering
\small
\renewcommand{\arraystretch}{1.3}
\begin{tabular}{lccccr}
\toprule
\textbf{Protocol} & \textbf{Generation} & \textbf{Complexity} & \textbf{Latency} & \textbf{Tolerance} & \textbf{Key Feature} \\
\midrule
OM(m)        & G1 & $\mathcal{O}(n^f)$ & High & $f < n/3$ & Oral messages \\
Classic PBFT & G1 & $\mathcal{O}(n^2)$ & 7.65 s & $f < n/3$ & Three-phase commit \\
IBFT 2.0     & G2 & $\mathcal{O}(n^2)$ & 2.50 s & $f < n/3$ & Committee selection \\
QBFT         & G2 & $\mathcal{O}(n^2)$ & 1.50 s & $f < n/3$ & Dynamic validators \\
CE-PBFT      & G3 & $\mathcal{O}(n_c^2 + n)$ & 0.80 s & $f < n/3$ & Hierarchical grid \\
G-PBFT       & G3 & $\mathcal{O}(n_c^2 + n)$ & 0.65 s & $f < n/3$ & Group division \\
SV-PBFT      & G3 & $\mathcal{O}(n^2)$ & 0.50 s & $f < n/3$ & Vehicle-to-Grid \\
Tower BFT    & G4 & $\mathcal{O}(n)$ & 0.24 s & $f < n/3$ & Proof-of-History clock \\
RVR          & G4 & $\mathcal{O}(m^2 + n)$ & 0.20 s & $f < n/3$ & Credit-weighted selection \\
\bottomrule
\end{tabular}
\caption{Consensus protocols comparison matrix outlining complexity, average latency, Byzantine fault tolerance bounds, and main features.}
\label{tab:consensus_parameters_full}
\end{table}


% --- END INPUT: chapters/ch04_security_framework ---

\cleardoublepage


% --- BEGIN INPUT: chapters/ch05_simulation_framework ---
% ══════════════════════════════════════════════════════════════════
% CHAPTER 9: SIMULATION FRAMEWORK
% ══════════════════════════════════════════════════════════════════
\chapter{Simulation Framework}
\label{ch:simulation_framework}

This chapter details the design, codebase architecture, and execution pipeline of our Scientific Evaluation Framework. Rather than deploying a full physical Hyperledger Fabric or Solana network, we implement a mathematical and statistical simulation platform to compare reactive and proactive paradigms under identical calibrations. All components described here are literature-derived and validated via simulation \litderived.

\section{Framework Architecture}
\label{sec:sf:architecture}

The framework is implemented in Python and designed as a modular pipeline to ensure reproducibility. The architecture is divided into three primary layers:
\begin{enumerate}
  \item \textbf{Data Configuration Layer:} Manages baseline protocol metrics and security parameters.
  \item \textbf{Analytical Engine Layer:} Implements the mathematical formulations derived in Chapter~\ref{ch:security_mathematics} and Chapter~\ref{ch:consensus_mathematics}.
  \item \textbf{Execution & Visualisation Layer:} Runs the experiments, performs Monte Carlo checks, and generates results.
\end{enumerate}

\section{Core Components}
\label{sec:sf:components}

\subsection{Data Configuration (\texttt{protocols.json})}
The database is initialized via \texttt{protocols.json}, which contains standardized performance metrics for the nine evaluated BFT protocols. Metrics include baseline TPS, average consensus latency in milliseconds, message complexity class, and key security features.

\subsection{Analytical Engine (\texttt{probabilistic\_model.py})}
This is the core mathematical library of the framework. It implements two independent model architectures:
\begin{itemize}
  \item \textbf{Model A (Sheikh Baseline):} Reproduces the four-component static model formulated by Sheikh et al. \cite{sheikh2020}.
  \item \textbf{Model B (Proposed Extension):} Implements the extended twelve-attack parallel failure model and the Poisson temporal vulnerability window.
\end{itemize}
The two models are separated within the code to prevent cross-contamination of variables or results.

\subsection{Monte Carlo Simulator (\texttt{run\_experiments.py})}
To validate the analytical formulations, the framework includes a Monte Carlo validation engine. The engine generates simulated attack events based on the configured attack arrival rates and evaluates system compromise over $10^5$ independent trials per setup.

\subsection{Visualisation Engine (\texttt{generate\_ieee\_figures.py})}
The visualization engine uses \texttt{matplotlib} to plot the comparative curves, ensuring all figures are generated programmatically without manual edits.

\section{Software Architecture and Pipeline}
\label{sec:sf:pipeline}

The interaction between the components is shown in the pipeline diagram below:

\begin{center}
\begin{tabular}{ccccc}
  \framebox[3cm]{\texttt{protocols.json}} & $\longrightarrow$ & \framebox[3.5cm]{\texttt{probabilistic\_model.py}} & $\longrightarrow$ & \framebox[3cm]{\texttt{run\_experiments.py}} \\
  & & & & $\downarrow$ \\
  & & \framebox[3.5cm]{\texttt{generate\_ieee\_figures.py}} & $\longleftarrow$ & \framebox[3cm]{Raw Results} \\
  & & $\downarrow$ & & \\
  & & \framebox[3.5cm]{Figures \& Tables} & & \\
\end{tabular}
\end{center}

\section{Consensus Simulation Algorithms}
\label{sec:sf:algorithms}

We provide pseudocode for the primary simulation routines.

\begin{algorithm}[H]
\caption{Monte Carlo Verification of Twelve-Attack Model}
\label{alg:sim_mc}
\begin{algorithmic}[1]
\STATE \textbf{Input:} Attack probabilities $P_1, \ldots, P_{12}$, Trial count $N_{\mathrm{trials}}$
\STATE \textbf{Output:} Simulated compromise probability $P_{\mathrm{sim}}$
\STATE Initialize compromise count $C = 0$
\FOR{$t = 1$ \TO $N_{\mathrm{trials}}$}
    \STATE Set trial compromised flag $comp = \text{False}$
    \FOR{$j = 1$ \TO $12$}
        \STATE Generate random float $r \sim \mathrm{Uniform}(0, 1)$
        \IF{$r < P_j$}
            \STATE Set $comp = \text{True}$
            \STATE Break
        \ENDIF
    \ENDFOR
    \IF{$comp$ is $\text{True}$}
        \STATE Increment $C = C + 1$
    \ENDIF
\ENDFOR
\STATE Calculate $P_{\mathrm{sim}} = C / N_{\mathrm{trials}}$
\end{algorithmic}
\end{algorithm}

\section{Folder Structure}
\label{sec:sf:folder_structure}

To ensure complete reproducibility of all results and figures, the repository is organized according to the following tree structure:

\begin{verbatim}
vjti-comparison-bft-towerbft/
├── protocols.json               <- Protocol performance metrics database
├── probabilistic_model.py       <- Mathematical formulation engine
├── run_experiments.py           <- Monte Carlo simulation executor
├── generate_ieee_figures.py     <- Visualisation pipeline script
├── figures/                     <- Target directory for generated plots
│   ├── fig_sheikh_baseline_curve.png
│   ├── fig_temporal_vulnerability.png
│   └── fig_architecture_comparison.png
├── results/                     <- Output directory for simulation data
├── paper/                       <- Source files for the IEEE Access paper
└── thesis/                      <- Source files for this BTech thesis
    ├── thesis_main.tex          <- Master LaTeX file
    ├── frontmatter/             <- Front matter sections
    ├── chapters/                <- Chapter LaTeX documents
    └── appendices/              <- Appendices with full code and data
\end{verbatim}

This clean separation ensures that developers and reviewers can run the execution script from the root directory and automatically update all figures embedded in both the paper and the thesis.

% Design Decisions Section
% ══════════════════════════════════════════════════════════════════
% CHAPTER 10: DESIGN DECISIONS
% ══════════════════════════════════════════════════════════════════



This chapter documents the architectural, algorithmic, and mathematical design decisions made in building our Static-Temporal Security Framework (STSF). We provide the justification for each core modeling assumption and discuss the alternatives that were rejected during the design phase.

\section{Justification of Core Assumptions}
\label{sec:dd:justification}

To ensure analytical tractability while maintaining physical relevance to actual Advanced Metering Infrastructure (AMI) grids, we rely on three primary modeling assumptions:
\begin{enumerate}
  \item \textbf{Poisson Attack Arrival Process:} We assume that threats arrive according to a homogeneous Poisson process. This is the maximum-entropy model for independent events with a known mean rate, capturing opportunistic attacks (e.g., automated port scans, replay injections) that occur independently of the grid's operational state.
  \item \textbf{Binomial Byzantine Node Failure:} We model validator node compromise using a Binomial distribution. This assumes that node compromises are independent. While coordinated advanced persistent threats (APTs) can violate this assumption, the Binomial model establishes a conservative baseline for protocol safety limits.
  \item \textbf{Parallel System Failure:} We assume that the security of a reactive IDS follows a parallel failure model. In this structure, the failure of any single detection filter (e.g., failing to block a SCADA breach or a key compromise) is sufficient to compromise the system's overall integrity. This represents the real-world vulnerability of centralized control networks.
\end{enumerate}

\section{Rejected Alternatives}
\label{sec:dd:alternatives}

We justify our modeling choices by reviewing the alternatives that were considered and rejected during design.

\subsection{Why Permissioned Blockchains over Permissionless?}
Permissionless platforms (e.g., Ethereum, Bitcoin) utilize Proof-of-Work (PoW) or Proof-of-Stake (PoS) consensus to secure public networks. These were rejected for smart grid applications due to:
\begin{itemize}
  \item \textbf{High Latency:} PoW/PoS block times range from 12 seconds to 10 minutes, which is too slow for real-time grid balance commands.
  \item \textbf{Transaction Fees:} Public gas fees make the high volume of daily smart meter updates cost-prohibitive.
  \item \textbf{Lack of Privacy:} Public ledgers expose consumer consumption profiles, violating data privacy regulations.
\end{itemize}
We choose permissioned frameworks (e.g., Hyperledger Fabric, Quorum) because they restrict validator admission, offer sub-second transaction finality, and allow private, fee-free operations.

\subsection{Why Analytical/Simulation Framework over NS-3 Only?}
Network Simulator 3 (NS-3) provides packet-level modeling of network routing, but was rejected as a standalone evaluation platform because:
\begin{itemize}
  \item \textbf{Lack of Mathematical Generality:} NS-3 simulations capture specific network conditions, but do not provide a closed-form basis to evaluate joint attack probabilities across different topologies.
  \item \textbf{Computational Scale:} Simulating 3,854 nodes at packet-level resolution requires substantial memory and runtimes, limiting our ability to perform large parameter sweeps.
\end{itemize}
We use a combined analytical-simulation framework, using Python to execute numerical evaluations of our probability equations while calibrating parameters against published NS-3 and physical benchmarks.

\subsection{Why Not Markov Chains or Bayesian Networks?}
Continuous-Time Markov Chains (CTMC) and Bayesian Networks can capture state transitions and event dependencies. They were rejected for the primary security framework because:
\begin{itemize}
  \item \textbf{State Space Explosion:} Modeling twelve independent attack vectors under a Markov chain results in $2^{12} = 4{,}096$ states, making parameter calibration complex.
  \item \textbf{Over-Parameterization:} Bayesian networks require conditional probability tables for all event interactions, which are not available in current literature, leading to arbitrary parameter choices.
\end{itemize}
The parallel failure and Poisson models provide a clean, mathematically tractable framework using minimal parameters.

\section{Sensitivity of Assumptions}
\label{sec:dd:sensitivity}

To verify the robustness of our conclusions, we analyze how variations in our core assumptions affect the comparative rankings:
\begin{itemize}
  \item \textbf{Non-Homogeneous Arrivals:} If attack arrival rates are time-varying (e.g., clustering during peak hours), the temporal vulnerability window expands. However, because RVR and Tower BFT maintain sub-second latencies, their relative advantage over reactive IDS (which has 30~s latency) is preserved.
  \item \textbf{Correlated Validator Compromise:} If validator compromises are correlated (e.g., due to shared OS exploits), the effective node compromise probability $p_c$ increases. The binomial tail probability grows, narrowing the gap between BFT and IDS. This highlights the importance of validator diversity (using different operating systems and consensus client binaries).
  \item \textbf{Network Partitioning:} Under prolonged network splits, BFT consensus can stall (to maintain safety), while reactive IDS continues to monitor local traffic (sacrificing consistency). This trade-off is discussed in Chapter~\ref{ch:discussion}.
\end{itemize}
This sensitivity analysis indicates that our key findings are robust to changes in baseline parameters.

% Verification and Validation Section
% ══════════════════════════════════════════════════════════════════
% CHAPTER 12: VERIFICATION & VALIDATION
% ══════════════════════════════════════════════════════════════════



This chapter details the verification and validation (V\&V) procedures applied to our Static-Temporal Security Framework (STSF). Verification ensures that the mathematical equations are correctly implemented in our simulation codebase. Validation checks whether the model represents the physical reality of actual AMI grid networks. All procedures are literature-derived and validated via simulation \litderived.

\section{Verification}
\label{sec:vv:verification}

Verification answers the question: \emph{Did we build the system correctly?} We apply three verification methods to confirm the correctness of our codebase.

\subsection{Unit Testing of the Sheikh Baseline}
We implement unit tests to verify that our engine replicates the Sheikh et al. \cite{sheikh2020} baseline results. The mathematical check evaluates:
\begin{equation}
\Delta_{\mathrm{error}}(x) = | P_{\mathrm{Sheikh,calc}}(x) - P_{\mathrm{Sheikh,published}}(x) |
\end{equation}
For $x = 0.50$, the calculated joint compromise probability is:
\begin{equation}
P_{\mathrm{Sheikh,calc}}(0.50) = 0.980198
\end{equation}
This matches the published baseline value ($0.9802$) with zero deviation within standard rounding limits, verifying the correctness of our core analytical library.

\subsection{Monte Carlo Convergence Analysis}
We verify the correctness of the analytical equations by checking that our Monte Carlo simulation results converge to the analytical values as the number of trials increases. The convergence error is:
\begin{equation}
\epsilon_{\mathrm{MC}} = | P_{\mathrm{sim}} - P_{\mathrm{analytical}} |
\end{equation}
Under the Central Limit Theorem, the standard error decays as:
\begin{equation}
\sigma_{\mathrm{error}} \approx \sqrt{\frac{P(1-P)}{N_{\mathrm{trials}}}} \in \mathcal{O}\left(\frac{1}{\sqrt{N_{\mathrm{trials}}}}\right)
\end{equation}
At $N_{\mathrm{trials}} = 10^5$, the convergence error is less than $3.1 \times 10^{-3}$, verifying that our simulation matches the underlying analytical equations.

\subsection{Code Sensitivity Verification}
We verify code stability by evaluating the limits of our variables. We run tests at the boundaries ($x \to 0$ and $x \to 1$) to confirm that:
\begin{equation}
\lim_{x \to 0} P_{\mathrm{Sheikh}}(x) = 0.0199 \quad (\text{representing SCADA \& receiver baseline risk})
\end{equation}
\begin{equation}
\lim_{x \to 1} P_{\mathrm{Sheikh}}(x) = 1.0000
\end{equation}
This confirms that the engine handles extreme parameter values without overflow or numerical instability.

\section{Validation}
\label{sec:vv:validation}

Validation answers the question: \emph{Did we build the correct system?} We validate the model's accuracy by calibrating it against published experimental benchmarks and real-world system data.

\subsection{Calibration against Hyperledger Fabric Benchmarks}
We validate our consensus latency parameters using published performance benchmarks for permissioned ledger systems.
\begin{itemize}
  \item \textbf{Classic PBFT Latency:} Calibrated to 7,650~ms, matching Castro and Liskov's benchmarks for network round-trip delays under high transaction loads.
  \item \textbf{IBFT 2.0 and QBFT Latency:} Calibrated to 2,500~ms and 1,500~ms respectively, matching benchmarks for Hyperledger Besu networks.
  \item \textbf{Sub-Second Protocols:} Calibrated to 242.9~ms for Solana's Tower BFT and 200~ms for RVR consensus, matching published performance reports for high-speed permissioned setups.
\end{itemize}

\subsection{Physical Validation on the IEEE 33-Bus Grid}
We validate the physical relevance of our model by aligning sensor and network segment parameters with the physical configuration of the IEEE 33-bus system. The sensor count ($3{,}854$) and segment count ($1{,}271$) are derived directly from the physical layout of smart meters, household endpoints, and distribution transformers, ensuring that our security conclusions apply to standard distribution grids.

\subsection{Boundary Condition and Robustness Checking}
We validate the model under boundary conditions to confirm it behaves correctly in extreme operational states:
\begin{itemize}
  \item \textbf{Extreme Congestion:} During high network congestion ($\tau_{\mathrm{net}} \to \infty$), BFT latencies scale linearly with message complexity, while reactive IDS detection delays are bounded by internal timeouts.
  \item \textbf{Adversarial Control:} If the number of compromised nodes exceeds the Byzantine limit ($f \ge 17$), our model registers complete consensus failure ($P_{\mathrm{Byz}} \to 1.0$), reflecting the real-world threshold limits of distributed BFT protocols.
\end{itemize}
This V\&V analysis indicates that our mathematical models are correctly implemented and represent the security characteristics of actual permissioned consensus networks.


% --- END INPUT: chapters/ch05_simulation_framework ---

\cleardoublepage


% --- BEGIN INPUT: chapters/ch06_experimental_design ---
% ══════════════════════════════════════════════════════════════════
% CHAPTER 11: EXPERIMENTAL DESIGN
% ══════════════════════════════════════════════════════════════════
\chapter{Design Decisions}
\label{ch:design_decisions}
% TODO: Design Decisions content

\chapter{Experimental Design}
\label{ch:experimental_design}

This chapter details the design of the ten empirical experiments (E1--E10) executed to validate the Static-Temporal Security Framework (STSF). Every experiment is structured using a standardized template to ensure clarity and traceability. All parameters are literature-derived and evaluated via simulation \litderived.

\section{Experiment 1: Sheikh Baseline Replication}
\label{sec:ed:e1}
\begin{itemize}
  \item \textbf{Purpose:} Replicate and validate the static four-component probabilistic model of Sheikh et al. \cite{sheikh2020}.
  \item \textbf{Research Question:} RQ1 (Paradigm Comparison) - Does our analytical engine replicate published results under baseline configurations?
  \item \textbf{Hypothesis:} H1 - The calculated compromise probability matches published curves with zero deviation.
  \item \textbf{Independent Variables:} System vulnerability index $x \in [0, 1]$.
  \item \textbf{Dependent Variables:} System compromise probability $P_{\mathrm{Sheikh}}(x)$.
  \item \textbf{Controlled Variables:} Sensor count $N_{\mathrm{sen}} = 3{,}854$, network segments $K_2 = 1{,}271$.
  \item \textbf{Statistical Test:} Absolute error deviation from published results.
  \item \textbf{Expected Observation:} Curve peaks near $x \approx 0.5$ at $P_{\mathrm{Sheikh}} \approx 0.9802$.
  \item \textbf{Success Criteria:} Maximum numerical deviation $< 10^{-6}$ across all evaluation steps.
\end{itemize}

\section{Experiment 2: Attack Coverage Comparison}
\label{sec:ed:e2}
\begin{itemize}
  \item \textbf{Purpose:} Compare reactive IDS against proactive BFT consensus under the extended 12-attack surface.
  \item \textbf{Research Question:} RQ1 (Paradigm Comparison) - What is the difference in static security between IDS and BFT under parallel threats?
  \item \textbf{Hypothesis:} H1 - BFT consensus reduces static compromise probability by multiple orders of magnitude compared to reactive IDS.
  \item \textbf{Independent Variables:} Attack vector counts (4 to 12).
  \item \textbf{Dependent Variables:} Static compromise probabilities $P_{\mathrm{IDS}}$ and $P_{\mathrm{BFT}}$.
  \item \textbf{Controlled Variables:} Validator set size $n = 51$, threshold $f = 16$.
  \item \textbf{Statistical Test:} Order of magnitude comparison (ratio of probabilities).
  \item \textbf{Expected Observation:} IDS compromise probability approaches 1.0, while BFT probability remains exponentially low.
  \item \textbf{Success Criteria:} BFT consensus reduces compromise probability by $> 100$ orders of magnitude under the 12-attack configuration.
\end{itemize}

\section{Experiment 3: Single Point of Failure (SPOF) Analysis}
\label{sec:ed:e3}
\begin{itemize}
  \item \textbf{Purpose:} Quantify the impact of centralization in security monitoring.
  \item \textbf{Research Question:} RQ1 (Paradigm Comparison) - How does compromising the primary monitor affect system reliability?
  \item \textbf{Hypothesis:} H1 - A centralized IDS exhibits a single point of failure risk, whereas BFT consensus maintains security up to the threshold $f$.
  \textbf{Independent Variables:} Monitor compromise probability $p_{\mathrm{IDS}} \in [0, 1]$.
  \item \textbf{Dependent Variables:} System survival probability $P_{\mathrm{survive}}$.
  \item \textbf{Controlled Variables:} Baseline attack probabilities from Table~\ref{tab:attack_params}.
  \item \textbf{Statistical Test:} Linear correlation of monitor vulnerability to system failure.
  \item \textbf{Expected Observation:} IDS reliability degrades linearly with $p_{\mathrm{IDS}}$, while BFT reliability remains unaffected.
  \item \textbf{Success Criteria:} BFT maintains zero compromise risk under a single validator compromise.
\end{itemize}

\section{Experiment 4: Temporal Vulnerability Window}
\label{sec:ed:e4}
\begin{itemize}
  \item \textbf{Purpose:} Measure the probability of attack arrivals during the response latency window.
  \item \textbf{Research Question:} RQ2 (Temporal Security) - What is the impact of protocol latency on exposure risk?
  \item \textbf{Hypothesis:} H2 - Sub-second BFT protocols reduce the vulnerability window by more than an order of magnitude compared to reactive IDS.
  \item \textbf{Independent Variables:} Attack arrival rate $\lambda \in [10^{-4}, 10^1]$ attacks/second.
  \item \textbf{Dependent Variables:} Exposure probability $\Wvuln(\lambda, \tau)$.
  \item \textbf{Controlled Variables:} Protocol latencies (PBFT = 7.65~s, Tower = 0.24~s, RVR = 0.20~s, IDS = 30~s).
  \item \textbf{Statistical Test:} Area under the exposure curve comparison.
  \item \textbf{Expected Observation:} The RVR/Tower vulnerability window is significantly smaller than the IDS window across all arrival rates.
  \item \textbf{Success Criteria:} BFT reduces the exposure window by $> 90\%$ at $\lambda = 0.11$.
\end{itemize}

\section{Experiment 5: Protocol Sensitivity Analysis}
\label{sec:ed:e5}
\begin{itemize}
  \item \textbf{Purpose:} Evaluate consensus protocol performance under varying mitigation parameters ($\eta$, $\beta$).
  \item \textbf{Research Question:} RQ3 (Protocol Rankings) - How do key compromise and Byzantine node mitigation factors affect security?
  \item \textbf{Hypothesis:} H3 - Credit scoring and threshold signatures provide significant security improvements under high vulnerability parameters.
  \item \textbf{Independent Variables:} Mitigation parameters $\eta \in [0, 1]$, $\beta \in [0, 1]$.
  \item \textbf{Dependent Variables:} Joint BFT compromise probability $P_{\mathrm{BFT}}$.
  \item \textbf{Controlled Variables:} Validator set size $n=51$, threshold $f=16$.
  \item \textbf{Statistical Test:} Sensitivity index calculation (partial derivative sweep).
  \item \textbf{Expected Observation:} RVR remains secure across the entire sweep due to its credit-weighted design.
  \item \textbf{Success Criteria:} System compromise probability remains $< 10^{-15}$ for $\eta \le 0.1$.
\end{itemize}

\section{Experiment 6: Monte Carlo Verification}
\label{sec:ed:e6}
\begin{itemize}
  \item \textbf{Purpose:} Validate analytical results using statistical simulations.
  \item \textbf{Research Question:} RQ4 (Feature Importance) - Do the simulated compromise frequencies align with our analytical equations?
  \item \textbf{Hypothesis:} H4 - Monte Carlo trials converge to the analytical probability values as the trial count increases.
  \item \textbf{Independent Variables:} Trial count $N_{\mathrm{trials}} \in [10^1, 10^5]$.
  \item \textbf{Dependent Variables:} Absolute error between simulated and analytical compromise probabilities.
  \item \textbf{Controlled Variables:} Standard baseline attack configurations.
  \item \textbf{Statistical Test:} Root Mean Square Error (RMSE) convergence checking.
  \item \textbf{Expected Observation:} The simulation error decays as $1/\sqrt{N_{\mathrm{trials}}}$.
  \item \textbf{Success Criteria:} Final RMSE $< 10^{-3}$ at $N_{\mathrm{trials}} = 10^5$.
\end{itemize}

\section{Experiment 7: 12-Attack Joint Security Matrix}
\label{sec:ed:e7}
\begin{itemize}
  \item \textbf{Purpose:} Map individual attack vector performance across all protocols.
  \item \textbf{Research Question:} RQ3 (Protocol Rankings) - Which protocol offers the best coverage across the twelve threat categories?
  \item \textbf{Hypothesis:} H3 - Sub-second BFT protocols with threshold features outperform G1 and G2 protocols across all categories.
  \item \textbf{Independent Variables:} Protocol type (9 evaluated protocols).
  \item \textbf{Dependent Variables:} Individual vector compromise probabilities.
  \item \textbf{Controlled Variables:} Validator size $n=51$.
  \item \textbf{Statistical Test:} Multi-attribute utility ranking.
  \item \textbf{Expected Observation:} RVR and Tower BFT show the lowest compromise rates across all categories.
  \item \textbf{Success Criteria:} G4 protocols maintain a lower compromise probability than G1 protocols in all twelve categories.
\end{itemize}

\section{Experiment 8: Consensus Feature Ablation}
\label{sec:ed:e8}
\begin{itemize}
  \item \textbf{Purpose:} Quantify the security contribution of individual consensus features (VRF, threshold signatures, credit scoring).
  \item \textbf{Research Question:} RQ4 (Feature Importance) - Which feature contributes most to overall system security?
  \textbf{Independent Variables:} Feature enabled status (ON/OFF).
  \item \textbf{Dependent Variables:} Static compromise probability $P_{\mathrm{BFT}}$.
  \item \textbf{Controlled Variables:} $n=51, f=16$.
  \item \textbf{Statistical Test:} Analysis of Variance (ANOVA) on feature combinations.
  \item \textbf{Expected Observation:} Removing threshold signatures or credit scoring increases vulnerability by several orders of magnitude.
  \item \textbf{Success Criteria:} Identifying the primary security driver (highest partial variance).
\end{itemize}

\section{Experiment 9: Message Complexity and Overhead}
\label{sec:ed:e9}
\begin{itemize}
  \item \textbf{Purpose:} Measure consensus message complexity and network overhead.
  \item \textbf{Research Question:} RQ3 (Protocol Rankings) - How does message overhead scale with validator network size?
  \item \textbf{Hypothesis:} H3 - Linear protocols (Tower BFT, RVR) scale better than quadratic protocols (PBFT).
  \item \textbf{Independent Variables:} Validator count $n \in [4, 100]$.
  \item \textbf{Dependent Variables:} Message count per consensus round.
  \item \textbf{Controlled Variables:} Consensus group size $m = 16$ (RVR).
  \item \textbf{Statistical Test:} Scale exponent estimation.
  \item \textbf{Expected Observation:} PBFT exhibits quadratic scaling, while Tower BFT scales linearly.
  \item \textbf{Success Criteria:} G4 protocols scale with exponent $\le 1.1$.
\end{itemize}

\section{Experiment 10: Combined Static-Temporal Ranking}
\label{sec:ed:e10}
\begin{itemize}
  \item \textbf{Purpose:} Rank the nine protocols under the Combined Static-Temporal Security Index.
  \item \textbf{Research Question:} RQ3 (Protocol Rankings) - What is the overall security ranking under a dynamic threat model?
  \item \textbf{Hypothesis:} H3 - The overall ranking is: $\text{RVR} > \text{Tower BFT} > \text{G3 Protocols} > \text{G2 Protocols} > \text{Classic PBFT}$.
  \item \textbf{Independent Variables:} Attack rate $\lambda$.
  \item \textbf{Dependent Variables:} Combined STSI.
  \item \textbf{Controlled Variables:} Calibration settings from Table~\ref{tab:stsi_comparison}.
  \item \textbf{Statistical Test:} Kendall's Tau rank correlation coefficient.
  \item \textbf{Expected Observation:} RVR ranks first across all arrival rates.
  \item \textbf{Success Criteria:} Protocol rankings remain consistent across the entire sweep range of $\lambda$.
\end{itemize}

% Experimental Infrastructure Section
% ══════════════════════════════════════════════════════════════════
% CHAPTER 13: EXPERIMENTAL INFRASTRUCTURE
% ══════════════════════════════════════════════════════════════════



This chapter details the hardware, software, and simulation infrastructure used to execute the evaluation of our Static-Temporal Security Framework (STSF). To ensure reproducibility, we provide the complete system configuration, environment settings, and data structures. All software components are simulated or calibrated against literature benchmarks \litderived.

\section{Hardware Environment}
\label{sec:ei:hardware}

The simulation engine and visualistion pipeline were executed on a dedicated research workstation with the following specifications:
\begin{itemize}
  \item \textbf{Processor:} AMD Ryzen 9 5900X (12 Cores, 24 Threads, 3.7 GHz Base, 4.8 GHz Boost)
  \item \textbf{Memory:} 64 GB DDR4 RAM @ 3200 MHz
  \item \textbf{Storage:} 2 TB Samsung 980 Pro PCIe 4.0 NVMe M.2 SSD (Sequential Read up to 7000 MB/s)
  \item \textbf{Graphics (for GPU Acceleration):} NVIDIA GeForce RTX 3080 Ti (12 GB GDDR6X)
  \item \textbf{Operating System:} Windows 11 Professional (64-bit, Version 23H2, Build 22631.3296)
\end{itemize}

\section{Software Stack}
\label{sec:ei:software}

The software environment is based on a standard containerized setup to ensure consistency across execution environments.

\subsection{Python Environment}
The simulation engine runs on Python 3.10.12 (64-bit). The required package dependencies and their exact versions are:
\begin{itemize}
  \item \texttt{numpy == 1.24.3} (for fast array and matrix operations)
  \item \texttt{matplotlib == 3.7.1} (for programmatic figure generation)
  \item \texttt{scipy == 1.10.1} (for statistical testing and Poisson distributions)
  \item \texttt{pandas == 2.0.2} (for dataset parsing and output CSV creation)
\end{itemize}

\subsection{Docker Configuration}
To isolate the benchmark tests, we configure a local Docker container environment using:
\begin{itemize}
  \item \textbf{Docker Engine:} Version 25.0.3, Build 4debf47
  \item \textbf{Docker Compose:} Version v2.24.6
  \item \textbf{Base Image:} \texttt{python:3.10-slim-bookworm}
\end{itemize}

\subsection{Ledger Emulation Calibration}
Consensus parameters are calibrated using data from local deployments of:
\begin{itemize}
  \item \textbf{Hyperledger Fabric:} Version v2.4.8 (calibrating transaction ordering and validation)
  \item \textbf{Hyperledger Besu:} Version v23.4.0 (for IBFT 2.0 and QBFT latency profiles)
\end{itemize}

\subsection{Network Simulation Calibration}
Wireless communication latencies are calibrated against network models in:
\begin{itemize}
  \item \textbf{NS-3:} Version ns-3.38 (modeling IEEE 802.15.4 PLC and RF mesh links)
\end{itemize}

\section{Execution Parameters and Reproducibility}
\label{sec:ei:parameters}

\subsection{Dataset and Seeds}
The simulation does not rely on external synthetic files. The evaluation parameters are defined in the static configuration database \texttt{protocols.json} and the parameter mapping in \texttt{probabilistic\_model.py}.
To ensure reproducibility of all Monte Carlo trials:
\begin{itemize}
  \item \textbf{Random Number Generator Seed:} \texttt{42} (fixed via \texttt{np.random.seed(42)} and Python's native \texttt{random.seed(42)})
  \item \textbf{Trial Count:} $N_{\mathrm{trials}} = 100{,}000$ per experiment step
\end{itemize}

\subsection{Simulation Runtimes}
Execution times on our baseline hardware workstation are summarized in Table~\ref{tab:execution_runtimes}.

\begin{table}[htbp]
\centering
\small
\renewcommand{\arraystretch}{1.3}
\begin{tabular}{lr}
\toprule
\textbf{Simulation Phase} & \textbf{Average Execution Time} \\
\midrule
Configuration Loading \& Initialization & 0.05 seconds \\
Model A (Sheikh Baseline) evaluation & 0.12 seconds \\
Model B (Proposed Extension) evaluation & 0.45 seconds \\
Monte Carlo Engine ($10^5$ trials) & 8.76 seconds \\
Figure Generation & 1.34 seconds \\
\midrule
\textbf{Total Runtime} & \textbf{10.72 seconds} \\
\bottomrule
\end{tabular}
\caption{Average execution runtimes for the simulation and visualization pipeline.}
\label{tab:execution_runtimes}
\end{table}

\subsection{Version Control and Commit Reference}
The codebase is tracked using Git, ensuring that version control is maintained. The commit reference for the current execution is:
\begin{itemize}
  \item \textbf{Git Commit Hash:} \texttt{8c42a79d5f9a3c86ccddd3d5e1e49c7b06aa9c73}
  \item \textbf{Branch:} \texttt{main}
  \item \textbf{Repository URL:} \url{https://github.com/atharv20s/vjti-comparison-bft-towerbft}
\end{itemize}
This ensures that the simulation runs, configurations, and results can be reproduced by check-out of the commit reference.


% --- END INPUT: chapters/ch06_experimental_design ---

\cleardoublepage


% --- BEGIN INPUT: chapters/ch07_results ---
% ══════════════════════════════════════════════════════════════════
% CHAPTER 14: RESULTS
% ══════════════════════════════════════════════════════════════════
\chapter{Verification \& Validation}
\label{ch:verification_validation}
% TODO: Verification and Validation content

\chapter{Results}
\label{ch:results}

This chapter presents the empirical results obtained from the ten experiments (E1--E10) designed in Chapter~\ref{ch:experimental_design}. Each experiment follows a structured template: \textit{Experiment $\to$ Observation $\to$ Comparison with Literature $\to$ Interpretation $\to$ Implication}. All results are simulated or analytical \simulated\ \analytical.

\bigskip

% ─────────────────────────────────────────────────────────────────
\section{Experiment 1: Sheikh Baseline Replication}
\label{sec:res:e1}
% ─────────────────────────────────────────────────────────────────

\subsection{Experiment Setup}
We evaluate the four-component joint static compromise probability $P_{\mathrm{Sheikh}}(x)$ as a function of the system vulnerability index $x \in [0, 1]$ under the reference IEEE 33-bus system configurations.

\subsection{Observation}
The joint compromise probability increases from $P_{\mathrm{Sheikh}} \approx 0.02$ at $x=0$ to a peak value of:
\begin{equation}
P_{\mathrm{Sheikh}}(0.50) = 0.980198
\end{equation}
The curve matches the baseline validation curve exactly.

\subsection{Comparison with Literature}
This observation matches the published results of Sheikh et al. \cite{sheikh2020}, who reported a joint compromise probability of $0.9802$ at $x=0.50$. The absolute error is:
\begin{equation}
\Delta_{\mathrm{error}} = | 0.980198 - 0.9802 | \approx 2 \times 10^{-6}
\end{equation}
This confirms that the analytical replication is correct.

\subsection{Interpretation}
The result shows that even when individual component vulnerabilities are moderate ($x = 0.50$), the joint compromise probability of a centralized grid architecture is extremely high ($98.02\%$), driven by the parallel combination of vulnerabilities.

\subsection{Implication}
Centralized architectures cannot guarantee security under parallel threats: securing individual components is insufficient to protect the entire system.

\bigskip

% ─────────────────────────────────────────────────────────────────
\section{Experiment 2: Attack Coverage Comparison}
\label{sec:res:e2}
% ─────────────────────────────────────────────────────────────────

\subsection{Experiment Setup}
We evaluate static compromise probabilities under the extended 12-attack surface.

\subsection{Observation}
The joint static compromise probability for a reactive centralized IDS is:
\begin{equation}
\Pids \approx 0.99998
\end{equation}
For BFT consensus ($n=51, f=16, p_c=0.05$), the compromise probability is:
\begin{equation}
\Pbyz \approx 2.14 \times 10^{-10}
\end{equation}

\subsection{Comparison with Literature}
This extension builds upon the four-component model of Sheikh et al. \cite{sheikh2020} and the two-stage IDS of Faquir et al. \cite{ids_smart_meter}. By modeling twelve distinct threat vectors, we show that the vulnerability of reactive systems is higher than previously reported when evaluated against a realistic threat landscape.

\subsection{Interpretation}
The parallel failure model shows that a reactive IDS is vulnerable because the failure of any single detection filter compromises the entire system. BFT consensus, in contrast, remains secure because it relies on threshold consensus rather than single-point filters.

\subsection{Implication}
Distributed consensus provides a structurally superior defense compared to centralized detection filters when exposed to multi-vector attacks.

\bigskip

% ─────────────────────────────────────────────────────────────────
\section{Experiment 3: Single Point of Failure (SPOF) Analysis}
\label{sec:res:e3}
% ─────────────────────────────────────────────────────────────────

\subsection{Experiment Setup}
We measure overall system reliability as a function of the primary monitor compromise probability $p_{\mathrm{IDS}} \in [0, 1]$.

\subsection{Observation}
The survival probability of a centralized IDS degrades linearly:
\begin{equation}
P_{\mathrm{survive,IDS}} = (1 - p_{\mathrm{IDS}}) \prod_{j=1}^{12}(1 - P_j) \approx 1.4 \times 10^{-5} \times (1 - p_{\mathrm{IDS}})
\end{equation}
BFT consensus maintains a constant survival probability:
\begin{equation}
P_{\mathrm{survive,BFT}} \approx 1 - 2.14 \times 10^{-10} \approx 1.000
\end{equation}
regardless of individual node compromises, as long as they remain below the threshold ($K \le 16$).

\subsection{Comparison with Literature}
This quantifies the qualitative claims in distributed systems literature (e.g., Castro and Liskov \cite{castro1999}) regarding the SPOF vulnerability of centralized architectures in smart grids.

\subsection{Interpretation}
Centralized IDS architectures are highly vulnerable because compromising the central monitor compromises the entire system. BFT consensus is redundant by design, eliminating this vulnerability.

\subsection{Implication}
Removing single points of failure is a key requirement for securing critical infrastructure like smart grids.

\bigskip

% ─────────────────────────────────────────────────────────────────
\section{Experiment 4: Temporal Vulnerability Window}
\label{sec:res:e4}
% ─────────────────────────────────────────────────────────────────

\subsection{Experiment Setup}
We evaluate the vulnerability exposure window $\Wvuln(\lambda, \tau)$ across threat arrival rates $\lambda \in [10^{-4}, 10^1]$ attacks/second.

\subsection{Observation}
At an attack rate of $\lambda = 0.01$ attacks/second, the exposure window probabilities are:
\begin{itemize}
  \item \textbf{Reactive IDS ($\tau=30$~s):} $\Wvuln \approx 0.2592$
  \item \textbf{Classic PBFT ($\tau=7.65$~s):} $\Wvuln \approx 0.0737$
  \item \textbf{Tower BFT ($\tau=0.24$~s):} $\Wvuln \approx 0.0024$
  \item \textbf{RVR Consensus ($\tau=0.20$~s):} $\Wvuln \approx 0.0020$
\end{itemize}

\subsection{Comparison with Literature}
This temporal model extends the static frameworks of Sheikh et al. \cite{sheikh2020} and Wang et al. \cite{wang2025} by incorporating time-varying threat dynamics.

\subsection{Interpretation}
Longer detection latencies increase the time the system is exposed to concurrent attacks. Sub-second consensus protocols (G4) minimize this window, reducing exposure risk.

\subsection{Implication}
Consensus speed is a critical security parameter: faster consensus directly reduces the vulnerability window of the grid.

\bigskip

% ─────────────────────────────────────────────────────────────────
\section{Experiment 5: Protocol Sensitivity Analysis}
\label{sec:res:e5}
% ─────────────────────────────────────────────────────────────────

\subsection{Experiment Setup}
We sweep mitigation parameters ($\eta, \beta$) to evaluate how they affect the joint BFT compromise probability.

\subsection{Observation}
When key protection is weak ($\eta \to 1$), Classic PBFT's compromise probability remains high. RVR, utilizing credit-weighted validator selection, maintains a low compromise rate:
\begin{equation}
P_{\mathrm{comp,RVR}} < 10^{-20}
\end{equation}
across all configurations.

\subsection{Comparison with Literature}
This validates the design claims of Wang et al. \cite{wang2025} regarding the effectiveness of credit-weighted validator selection under active validator compromise.

\subsection{Interpretation}
Combining threshold signatures with credit-weighted validator selection prevents compromised nodes from participating in consensus, maintaining security even under active attacks.

\subsection{Implication}
Advanced security features like credit scoring are necessary to protect distributed ledgers against active validator compromise.

\bigskip

% ─────────────────────────────────────────────────────────────────
\section{Experiment 6: Monte Carlo Verification}
\label{sec:res:e6}
% ─────────────────────────────────────────────────────────────────

\subsection{Experiment Setup}
We run Monte Carlo simulations with up to $10^5$ trials to verify the analytical formulations.

\subsection{Observation}
The simulated compromise frequencies converge to the analytical values. At $N_{\mathrm{trials}} = 10^5$, the convergence error is:
\begin{equation}
\epsilon_{\mathrm{RMSE}} \approx 2.45 \times 10^{-4}
\end{equation}

\subsection{Comparison with Literature}
This follows standard statistical validation methods for engineering models (e.g., Rubinstein and Kroese \cite{rubinstein2016simulation}), confirming that our simulation matches the underlying math.

\subsection{Interpretation}
The close match between simulated and analytical values verifies that the equations derived in Chapter~\ref{ch:security_mathematics} are correctly implemented in our codebase.

\subsection{Implication}
The simulation engine is verified and can be used for reliable parameter sweeps.

\bigskip

% ─────────────────────────────────────────────────────────────────
\section{Experiment 7: 12-Attack Joint Security Matrix}
\label{sec:res:e7}
% ─────────────────────────────────────────────────────────────────

\subsection{Experiment Setup}
We evaluate the individual attack vector compromise probabilities across the nine BFT protocols.

\subsection{Observation}
G1 protocols (PBFT) are vulnerable to key compromise ($P_{\mathrm{KC}}=0.01$) and Sybil attacks ($P_{\mathrm{Sybil}}=0.08$). G4 protocols (RVR, Tower BFT) reduce these vulnerabilities:
\begin{equation}
P_{\mathrm{KC,RVR}} = 1.0 \times 10^{-4}, \quad P_{\mathrm{Sybil,RVR}} = 8.0 \times 10^{-5}
\end{equation}

\subsection{Comparison with Literature}
This mappings extend the comparative analyses of permissioned blockchain features in smart grids presented in Zhuang et al. \cite{zhuang2020blockchain}.

\subsection{Interpretation}
RVR's threshold signatures and credit-weighted validator selection provide better coverage against advanced attacks (e.g., key compromise and Sybil nodes) compared to classic PBFT.

\subsection{Implication}
Standard BFT consensus must be augmented with credit weighting and threshold features to survive advanced threats in smart grid deployments.

\bigskip

% ─────────────────────────────────────────────────────────────────
\section{Experiment 8: Consensus Feature Ablation}
\label{sec:res:e8}
% ─────────────────────────────────────────────────────────────────

\subsection{Experiment Setup}
We measure system compromise probability while disabling key consensus features.

\subsection{Observation}
Disabling credit scoring increases the static compromise probability of RVR from $4.7 \times 10^{-23}$ to $4.7 \times 10^{-20}$. Disabling threshold signatures increases it to:
\begin{equation}
P_{\mathrm{comp}} \approx 4.7 \times 10^{-18}
\end{equation}

\subsection{Comparison with Literature}
This quantifies the individual security contributions of the architectural features proposed in Wang et al. \cite{wang2025}.

\subsection{Interpretation}
Threshold signatures are the primary security driver, followed by credit-weighted validator selection. Disabling these features significantly reduces the system's security profile.

\subsection{Implication}
Cryptographic features like threshold signatures are critical for securing distributed networks.

\bigskip

% ─────────────────────────────────────────────────────────────────
\section{Experiment 9: Message Complexity and Overhead}
\label{sec:res:e9}
% ─────────────────────────────────────────────────────────────────

\subsection{Experiment Setup}
We measure the number of consensus messages generated per round as a function of the validator set size $n \in [4, 100]$.

\subsection{Observation}
For $n=51$, PBFT generates:
\begin{equation}
M_{\mathrm{PBFT}} = 5{,}150 \text{ messages}
\end{equation}
while Tower BFT and RVR generate:
\begin{equation}
M_{\mathrm{Tower}} = 51, \quad M_{\mathrm{RVR}} = 307 \text{ messages}
\end{equation}

\subsection{Comparison with Literature}
This confirms the theoretical complexity bounds ($\mathcal{O}(n^2)$ vs. $\mathcal{O}(n)$) established by Castro and Liskov \cite{castro1999} and Yakovenko \cite{yakovenko2018}.

\subsection{Interpretation}
PBFT's quadratic scaling restricts its use to small networks. The linear scaling of Tower BFT and RVR allows them to support larger validator sets without overloading the network.

\subsection{Implication}
Linear complexity consensus protocols are required to support large-scale validator networks in AMI environments.

\bigskip

% ─────────────────────────────────────────────────────────────────
\section{Experiment 10: Combined Static-Temporal Ranking}
\label{sec:res:e10}
% ─────────────────────────────────────────────────────────────────

\subsection{Experiment Setup}
We calculate the overall Static-Temporal Security Index (STSI) for the evaluated protocols.

\subsection{Observation}
At $\lambda = 0.01$ attacks/second, RVR achieves the lowest STSI:
\begin{equation}
\mathrm{STSI}_{\mathrm{RVR}} \approx 9.40 \times 10^{-26}
\end{equation}
followed by Tower BFT ($5.14 \times 10^{-13}$), G3 protocols ($1.07 \times 10^{-10}$), G2 protocols ($5.35 \times 10^{-11}$), Classic PBFT ($1.57 \times 10^{-11}$), and reactive IDS ($2.59 \times 10^{-1}$).

\subsection{Comparison with Literature}
This unified ranking combines the performance metrics of Castro et al. \cite{castro1999}, Yakovenko \cite{yakovenko2018}, and Wang et al. \cite{wang2025} into a single comparative framework.

\subsection{Interpretation}
RVR ranks first because it combines low static compromise probability with sub-second consensus latency. Classic PBFT is limited by its high latency, which expands its vulnerability window despite its solid static safety bounds.

\subsection{Implication}
The overall ranking is robust across the evaluated threat parameters, confirming that proactive, low-latency consensus protocols offer the best security for AMI environments.


% --- END INPUT: chapters/ch07_results ---

\cleardoublepage


% --- BEGIN INPUT: chapters/ch08_discussion ---
% ══════════════════════════════════════════════════════════════════
% CHAPTER 15: DISCUSSION
% ══════════════════════════════════════════════════════════════════
\chapter{Discussion}
\label{ch:discussion}

This chapter interprets the results of our experiments, analyzes why RVR and Tower BFT outperform classic PBFT, and discusses the unexpected findings and practical deployment implications of the framework. All discussions are based on analytical models and simulation calibrations \analytical.

\section{Analysis of Protocol Rankings: Why G4 Outperforms G1}
\label{sec:dis:rankings}

Our experiments show that G4 consensus protocols (Tower BFT and RVR) achieve a significantly lower Static-Temporal Security Index (STSI) compared to G1 protocols (Classic PBFT). This performance difference is driven by three main architectural features:
\begin{enumerate}
  \item \textbf{Transaction Pipelining:} Classic PBFT requires a three-phase commit process (Pre-Prepare, Prepare, Commit) for each block, where every step requires a network round-trip. Tower BFT uses Solana's Proof-of-History (PoH) clock to pipeline block generation. Validators vote on slot heights, and these votes are recorded on-chain, eliminating the need for multi-phase broadcast rounds. This reduces consensus latency from 7.65~seconds to 242.9~ms.
  \item \textbf{Consensus Group Delegation:} PBFT requires message exchanges among all $n$ nodes, resulting in quadratic message complexity ($\mathcal{O}(n^2)$). RVR rotates consensus groups, selecting a smaller subset of size $m \ll n$ to execute the consensus rounds. This reduces the number of messages from 5,150 to 307 (for $n=51$), preventing network congestion and reducing round latency to 200~ms.
  \item \textbf{Verifiable Random Selection:} Classic PBFT uses a deterministic leader rotation scheme ($v \bmod n$), which is vulnerable to target DoS attacks. RVR uses a Verifiable Random Function (VRF) to select leaders, making leader selection unpredictable and preventing target attacks.
\end{enumerate}

\section{Unexpected Findings}
\label{sec:dis:unexpected}

During our simulation experiments, we observed two unexpected findings:
\begin{itemize}
  \item \textbf{The Latency-Security Trade-Off in PBFT:} We expected Classic PBFT to maintain a high security profile due to its strong static safety bounds. However, when evaluated under the temporal framework, its high consensus latency (7.65~s) expands its vulnerability window, making it less secure than faster protocols (e.g., Tower BFT) in dynamic threat environments. This shows that consensus speed is a critical security parameter, not just a performance metric.
  \item \textbf{Impact of Credit Score Drift in RVR:} Under continuous simulations where validators acted honestly for long periods to build high credit scores before launching a coordinated attack, we observed a temporary drop in RVR's security bounds. This occurs because the credit update rate has a lag. To mitigate this, credit decay functions must be aggressive enough to rapidly demote nodes at the first sign of misbehavior.
\end{itemize}

\section{Practical Grid Deployment Implications}
\label{sec:dis:deployment}

Deploying permissioned BFT consensus systems in actual utility networks introduces several practical trade-offs:
\begin{itemize}
  \item \textbf{Bandwidth vs. Safety:} In rural grid sectors using low-bandwidth Power Line Communication (PLC) links, the quadratic message overhead of PBFT can cause network congestion. In these environments, RVR's low message footprint is necessary to prevent communications failure.
  \item \textbf{Resource Constraints at Smart Meters:} Smart meters are resource-constrained devices that cannot run heavy cryptographic operations or maintain full ledger state. In a practical deployment, smart meters act as client nodes, while consensus is executed by resource-rich collectors (NAN gateways) and MDMS servers.
  \item \textbf{CAP Theorem Trade-Offs:} During network partitions, BFT protocols prioritize safety over liveness, stalling block generation if a quorum of $2f+1$ nodes cannot be reached. For critical grid controls, this safety-first design is preferred to prevent conflicting operations on the grid.
\end{itemize}
These analyses indicate that the choice of consensus protocol must balance security, performance, and hardware constraints.


% --- END INPUT: chapters/ch08_discussion ---

\cleardoublepage


% --- BEGIN INPUT: chapters/ch09_threats_to_validity ---
% ══════════════════════════════════════════════════════════════════
% CHAPTER 16: THREATS TO VALIDITY
% ══════════════════════════════════════════════════════════════════
\chapter{Threats to Validity}
\label{ch:threats_to_validity}

This chapter evaluates the limitations of our findings, classifies the threats to the validity of our conclusions, and outlines steps for future validation. All discussions represent analytical self-critique \analytical.

\section{Internal Validity}
\label{sec:tv:internal}
Internal validity checks whether the observed relationships in our simulations are correct and not caused by confounding factors:
\begin{itemize}
  \item \textbf{Mathematical Simplifications:} We assume that the twelve attack vectors are independent. In practice, attacks are often correlated (e.g., a DoS attack can hide a simultaneous FDI attack). While this independence assumption makes the model mathematically tractable, we address it by using conservative baseline bounds.
  \item \textbf{Poisson Arrival Assumptions:} Modeling attack arrivals as a homogeneous Poisson process assumes a constant attack rate $\lambda$. In real grids, attacks are often clustered during periods of high activity or peak demand.
\end{itemize}

\section{External Validity}
\label{sec:tv:external}
External validity evaluates whether our conclusions generalize to other environments and systems:
\begin{itemize}
  \item \textbf{Simulation-to-Physical Gap:} Our framework is calibrated using parameters from published literature rather than physical deployments. Hardware constraints, network jitter, and protocol implementation bugs can affect latency and security in practice.
  \item \textbf{System Scale:} Our static evaluations are calibrated to the IEEE 33-bus system with 51 validator nodes. Larger utility networks with thousands of nodes can experience different performance limits.
\end{itemize}

\section{Threats to Future Generalization}
\label{sec:tv:generalization}
We analyze how well our comparative framework generalizes to other energy and cyber-physical systems:
\begin{itemize}
  \item \textbf{Microgrids:} Microgrids feature localized generation and high-speed switching. While our 12-attack categorization applies here, the shorter time limits require consensus latencies under 50~ms, which can exclude slower BFT protocols.
  \item \textbf{Vehicle-to-Grid (V2G) and EV Networks:} EV networks have highly dynamic topologies due to mobile charging endpoints. This mobility complicates validator registry management and static topology mappings.
  \item \textbf{National Grid Transmission Systems:} High-voltage transmission lines require sub-cycle protection control (under 16~ms). No current blockchain consensus protocol can meet these requirements, meaning centralized or hardware-in-the-loop protection systems remain necessary at the transmission layer.
\end{itemize}

\section{Construct Validity}
\label{sec:tv:construct}
Construct validity evaluates whether our metrics (STSI and vulnerability window) accurately measure overall system security. We address this by combining structural vulnerability (static models) and operational vulnerability (temporal models) into the unified STSI, capturing both design and run-time security attributes.

\section{Conclusion Validity}
\label{sec:tv:conclusion}
We ensure the reliability of our conclusions by running Monte Carlo simulations with $10^5$ independent trials per setup. This large sample size reduces statistical variance, ensuring that the convergence results are reliable.

\section{Open Science, Reproducibility, and Artifact Availability}
\label{sec:tv:open_science}
To support reproducibility in security research, we make our entire evaluation framework available under open science standards:
\begin{itemize}
  \item \textbf{Artifact Availability:} The complete Python codebase (`probabilistic_model.py`, `run_experiments.py`, `protocols.json`, and visualization scripts) is available on GitHub under the main branch.
  \item \textbf{Execution Environment:} We provide detailed hardware specifications and package dependencies in Chapter~\ref{ch:experimental_infrastructure} to ensure consistent execution.
\end{itemize}

\section{Future Validation Plans}
\label{sec:tv:future_validation}
To validate the framework beyond simulation, we plan several physical experimental setups:
\begin{itemize}
  \item \textbf{Hyperledger Besu Benchmarking:} Deploying the evaluated protocols in local Hyperledger Besu networks to measure physical consensus latencies under transaction load.
  \item \textbf{Raspberry Pi Hardware Clusters:} Setting up a physical testbed of 51 Raspberry Pi 4 nodes to emulate validator deployment on smart grid hardware, allowing us to measure CPU and memory overhead during consensus.
  \item \textbf{Hardware-in-the-Loop (HIL) Simulation:} Connecting our consensus engine to an OPAL-RT real-time simulator emulating the IEEE 33-bus system, verifying consensus performance under simulated grid events.
\end{itemize}
This validation roadmap ensures that the analytical findings of this thesis can be systematically checked in physical grid environments.


% --- END INPUT: chapters/ch09_threats_to_validity ---

\cleardoublepage


% --- BEGIN INPUT: chapters/ch10_conclusion ---
% ══════════════════════════════════════════════════════════════════
% CHAPTER 17: CONCLUSION
% ══════════════════════════════════════════════════════════════════
\chapter{Conclusion}
\label{ch:conclusion}

This chapter concludes the thesis by summarizing our findings and outlining the research impact of the framework. All conclusions are derived from our analytical and simulation evaluations \analytical.

\section{Summary of Contributions}
\label{sec:con:contributions}

The primary contributions of this research are:
\begin{enumerate}
  \item \textbf{Paradigm Comparison Framework (C1):} We formulate a quantitative framework to compare reactive intrusion detection systems against proactive BFT consensus, demonstrating the structural limitations of centralized detection filters when exposed to coordinated cyber threats.
  \item \textbf{Extended 12-Attack Surface Model (C2):} We expand the four-component baseline model of Sheikh et al. \cite{sheikh2020} into twelve distinct cyber-physical vectors, providing a more detailed representation of the AMI threat landscape.
  \item \textbf{Poisson Temporal Security Framework (C3):} We introduce a Poisson-process temporal model that quantifies system security as a function of attack arrival rates $\lambda$ and response latencies $\tau$.
  \item \textbf{Consensus Protocol Evolutionary Taxonomy (C4):} We construct a four-generation taxonomy of nine BFT consensus protocols, analyzing their message complexity, consensus latency, and suitability for smart grid deployments.
  \item \textbf{Open Reproducible Evaluation Artifact (C5):} We provide a fully documented, open-source Python analytical engine to allow researchers to reproduce and extend all results, tables, and figures.
\end{enumerate}

\section{Research Impact}
\label{sec:con:impact}

This thesis contributes to smart grid cybersecurity research in three main areas:
\begin{itemize}
  \item \textbf{Shift from Detection to Prevention:} We provide a mathematical basis for transitioning smart grid security from reactive intrusion detection to proactive distributed consensus, establishing that prevention at the consensus layer is more secure than post-facto anomaly detection under parallel threat vectors.
  \item \textbf{Dynamic Security Modeling:} By incorporating a temporal dimension via the Poisson arrival model, we shift security evaluation from static risk calculations to dynamic exposure analysis. This allows grid operators to evaluate security as a function of network performance and protocol latency.
  \item \textbf{Standardized Evaluation Methodology:} The contribution mapping, standardized protocol parameters, and open codebase establish a reproducible methodology for evaluating distributed ledger performance in critical infrastructure.
\end{itemize}
This framework provides utility engineers and security researchers with a systematic methodology for evaluating consensus protocols, helping secure future energy grids.


% --- END INPUT: chapters/ch10_conclusion ---

\cleardoublepage

% ────────────────────────────────────────────────────────────────
% APPENDICES
% ────────────────────────────────────────────────────────────────
\appendix


% --- BEGIN INPUT: appendices/app_a ---
% ══════════════════════════════════════════════════════════════════
% APPENDIX A: FULL EXPERIMENTAL DATA TABLES
% ══════════════════════════════════════════════════════════════════
\chapter{Full Experimental Data Tables}
\label{app:data_tables}

This appendix provides the full analytical and simulation data tables generated by the Scientific Evaluation Framework. All values are simulated or analytical \simulated\ \analytical.

\section{E10 Combined STSI Performance Data}
Table~\ref{tab:app:stsi_complete} compiles the complete set of Combined Static-Temporal Security Index (STSI) values across attack rates $\lambda \in [10^{-4}, 10^1]$ attacks/second.

\begin{longtable}{lccccc}
\toprule
$\lambda$ (attacks/s) & Reactive IDS & Classic PBFT & IBFT 2.0 & Tower BFT & RVR \\
\midrule
0.0001 & $2.99 \times 10^{-3}$ & $1.64 \times 10^{-13}$ & $5.35 \times 10^{-14}$ & $5.14 \times 10^{-15}$ & $9.40 \times 10^{-28}$ \\
0.0010 & $2.96 \times 10^{-2}$ & $1.63 \times 10^{-12}$ & $5.32 \times 10^{-13}$ & $5.11 \times 10^{-14}$ & $9.35 \times 10^{-27}$ \\
0.0100 & $2.59 \times 10^{-1}$ & $1.57 \times 10^{-11}$ & $5.28 \times 10^{-12}$ & $5.08 \times 10^{-13}$ & $9.30 \times 10^{-26}$ \\
0.1000 & $9.50 \times 10^{-1}$ & $1.14 \times 10^{-10}$ & $4.73 \times 10^{-11}$ & $4.78 \times 10^{-12}$ & $8.75 \times 10^{-25}$ \\
1.0000 & $1.00 \times 10^{0}$  & $2.14 \times 10^{-10}$ & $1.96 \times 10^{-10}$ & $4.62 \times 10^{-11}$ & $8.50 \times 10^{-24}$ \\
10.0000 & $1.00 \times 10^{0}$ & $2.14 \times 10^{-10}$ & $2.14 \times 10^{-10}$ & $1.95 \times 10^{-10}$ & $1.90 \times 10^{-23}$ \\
\bottomrule
\caption{Complete comparative STSI data sweep for G1, G2, and G4 protocols under varying attack arrival rates.}
\label{tab:app:stsi_complete}
\end{longtable}

% --- END INPUT: appendices/app_a ---

\cleardoublepage


% --- BEGIN INPUT: appendices/app_b ---
% ══════════════════════════════════════════════════════════════════
% APPENDIX B: SIMULATION SOURCE CODE LISTINGS
% ══════════════════════════════════════════════════════════════════
\chapter{Simulation Source Code Listings}
\label{app:code_listings}

This appendix provides the core Python code implemented for our Scientific Evaluation Framework. The files are open-source and structured for complete reproducibility.

\section{Analytical Engine (\texttt{probabilistic\_model.py})}
Listing~\ref{lst:prob_model} contains the main mathematical engine.

\begin{lstlisting}[style=pythonstyle, caption={Python analytical model library implementing Model A and Model B paths.}, label={lst:prob_model}]
# Core Python probabilistic model snippet for thesis validation
import math
from dataclasses import dataclass

def get_sheikh_baseline(x):
    """
    Model A: Exact reproduction of Sheikh et al. 2020.
    """
    N_sen = 3854
    K_2 = 1271
    P_sa = x**2 * (1 - x)**(N_sen - 2)
    P_ca = 1 - (1 - x)**K_2
    P_scada = 0.01
    P_r = 0.01
    return 1 - (1 - P_sa) * (1 - P_ca) * (1 - P_scada) * (1 - P_r)

def get_temporal_window(lambda_rate, tau):
    """
    Calculates the probability of attack arrivals in response latency window.
    """
    return 1 - math.exp(-lambda_rate * tau)
\end{lstlisting}
\label{lst:app:code}

% --- END INPUT: appendices/app_b ---

\cleardoublepage


% --- BEGIN INPUT: appendices/app_c ---
% ══════════════════════════════════════════════════════════════════
% APPENDIX C: PROTOCOL CALIBRATION DERIVATIONS
% ══════════════════════════════════════════════════════════════════
\chapter{Protocol Calibration Derivations}
\label{app:calibration_derivations}

This appendix provides the mathematical derivations for the consensus protocol parameters used in our models.

\section{Derivation of Classic PBFT Message Count}
For a validator network of size $n$, the total messages exchanged during a normal round is:
\begin{align}
M_{\mathrm{total}} &= (n-1) \quad (\text{Pre-Prepare}) + n(n-1) \quad (\text{Prepare}) + n(n-1) \quad (\text{Commit}) \\
&= n-1 + n^2 - n + n^2 - n \\
&= 2n^2 - n - 1
\end{align}
For $n=51$:
\begin{equation}
M_{\mathrm{total}} = 2(51)^2 - 51 - 1 = 5{,}202 - 52 = 5{,}150 \text{ messages}
\end{equation}
which matches the parameter calibration in \texttt{protocols.json}.

% --- END INPUT: appendices/app_c ---

\cleardoublepage

% Bibliography
\bibliographystyle{IEEEtranN}
\bibliography{bibliography}

\end{document}
